\begin{theorem}[Source result P3K.114 under continuation review]\label{night:P3K.114}
\StatementRecord{night:P3K.114}{open}{}{\cref{night:P3K.96}, \cref{night:P3K.109}, \cref{night:P3K.110}, \cref{night:P3K.113}; the original source, point line, roots, characters, reference, selected fan and scalar range stated here}
Retain the original marked type-III source
\[
 E_0:\ y^2=x^3+tx+t^2,\qquad G=(0,t),\qquad
 t_c=-4/27,\quad q=t+4/27,
\]
the indicated simple selected node with its actual branch marking,
and any prescribed lattice-basis fan. Use the genuine point line
\(M_a=\mathcal O([aG]-[O])\), chronological source translation by
\(-bG\), and either complete infinity-cap progression
\[
 (a,b;m)=(6,1;3l+2),\ l\ge0,\qquad
 (a,b;m)=(2,3;3l),\ l\ge1.
\]
Keep both primitive norms \(n=\pm1\), every prescribed finite
character \(\zeta_m^d\), the whole P3K.96 annular reference, and a
nonzero original scalar in the common strict contraction domain.
At the finite III fibre retain the inherited order-one, norm-zero
filling of the P3K.96 family.
Let \(X\) be the resulting completed threefold.

Every holomorphic automorphism fixes the original base pointwise,
and multiplication on the actual point line induces the exact
group isomorphism
\[
       \mathbb C^*_{\rm physical}\xrightarrow{\ \cong\ }
                       \operatorname{Aut}_{\rm hol}(X).
\]
For the full scalar family with \(\lambda=R^m\), every holomorphic
endpoint identification has integral scalar winding, and its
complete winding group and primitive forward maps are
\[
\begin{array}{c|c|c|c}
 (a,b;m)&\text{winding group}&\text{ordinary map}&\text{cap map}\\ \hline
 (6,1;3l+2)&m\mathbb Z&q^n&h^n\\
 (2,3;3l)&l\mathbb Z&(R^{-l}tU_G)^n&(R^{-l}U'_G)^n ,
\end{array}
\]
where \(h=1+4z^m/27=q_\infty q\),
\(U_G(p,v)=(p+G,B_{a,1}(p+G)^{-1}v)\), and
\(U'_G\) is its full good-root counterpart. The point coefficient
is \(B_{a,1}(p)=f_2(p)^a f_a(p-G)/f_a(p)\), with the full
Miller and Cartier arrows specified below. Their actual inverses
are the inverse-endpoint return maps.

The three chronological primitive M2 endpoints have composite
\(q^nB_\lambda^{-n}\) on the ordinary cover and
\(h^nB_\lambda^{-n}\) on the root cover; thus they give precisely
the canonical \(q^n/h^n\) pair on the quotient. Every negative
linear endpoint is excluded by the complete annular deck
\(B_\lambda^{n+m\nu}\), retaining its arbitrary reference integer
\(\nu\). The ratio of any two endpoints with the same winding
is exactly a physical scalar. All source, root, selected-fan,
Cartier and finite-character maps in these assertions are the
displayed original maps. No numerical return is assigned.
\end{theorem}
\begin{proof}[Supplied source argument; continuation review in progress]

% Expanded verified cp36 input: notes/part3_e6_iii_fixed_source_endpoint_proof.tex
This separate continuation keeps K109, K110 and EB unchanged. The source
is the actual \(r=0\) member of the K96 disc family,

\[
 E_0:\quad y^2=x^3+tx+t^2,\qquad G=(0,t),\qquad
 t_c=-4/27,\quad q=t+4/27.
 \tag{R0E.1}
\]

\par

\medskip\noindent\textbf{Retained finite-stage history.}

\par

The following scope paragraph and R0E.1-\/-3 record the completed initial finite-chart stage. Their former global-proof obligation is resolved by the full R0E.4-\/-10 proof reproduced below. The historical scope paragraph is retained verbatim in content.

The following calculation concerns the finite III source and its point
line \(M_2=\mathcal O(2G-O)\). It records the full coefficient before
any central restriction. The completed global period and endpoint
classification will be written separately after these charts are
checked; this note presently assigns no automorphism group or return.

\par

\medskip\noindent\textbf{R0E.1. The genuine quartic source cover and its line coefficient}

\par\nobreak

Use the actual map

\[
 t=z^4,\quad x=z^2X,\quad y=z^3Y,\qquad
 Y^2=X^3+X+z^2,\quad G'=(0,z).
 \tag{R0E.2}
\]

Its deck is \(D(z,X,Y)=(iz,-X,iY)\). The differential is

\[
 \omega_t=\frac{dx}{2y}=z^{-1}\frac{dX}{2Y}
 \tag{R0E.3}
\]

on each fibre. Doubling the actual section gives

\[
 2G=(1/4,-t-1/8),\qquad
 \tau(2G')=-\frac{X(2G')}{Y(2G')}
            =\frac{2z}{1+8z^4}.
 \tag{R0E.4}
\]

The finite point-to-tensor coefficient and its entire pullback are

\[
 f_2=\frac{y-t-x/2}{x-1/4}=z^2\phi,
 \qquad
 \phi=\frac{zY-z^2-X/2}{z^2X-1/4}.
 \tag{R0E.5}
\]

For the actual translation \(T:p\mapsto p-G'\), set

\[
 L=\frac{Y+z}{X},\qquad
 X'=L^2-X=\frac1X+\frac{2z(Y+z)}{X^2},\qquad
 Y'=L(X-X')-Y.
 \tag{R0E.6}
\]

These are rational expressions for the source map, with their actual
Cartier-module extensions. They preserve the whole equation R0E.2.
The entire coefficient of the point-line lift is

\[
 \frac{B_{2,1}}{t}
     =\phi(X,Y,z)\,\phi(X',Y',z).
 \tag{R0E.7}
\]

At \(z=0\), the source arrow is exactly

\[
 (X,Y)\longmapsto(X^{-1},-Y/X^2),\qquad
 \phi|_{z=0}=2X,
 \qquad (B_{2,1}/t)|_{z=0}=4.
 \tag{R0E.8}
\]

The equality is in the Hom line; the zeros and poles of \(2X\)
are retained by the two point/O Cartier modules. The original-source
components require the graph charts below, rather than a replacement
of those modules by the number \(4\).

\par

\medskip\noindent\textbf{R0E.2. Both original III components and the actual point frames}

\par\nobreak

On the strict cuspidal component use \(p=-x/y\), so

\[
 x=p^{-2},\quad y=-p^{-3}\qquad(t=0).
 \tag{R0E.9}
\]

On the other component use the literal resolution chart

\[
 x=t\alpha,\quad y=t\beta,\qquad
 \beta^2=1+\alpha+t\alpha^3.
 \tag{R0E.10}
\]

The two source arrows induced by \(p\mapsto p-G\), before the
point-line factor is read, are

\[
 C_0\longrightarrow C_1:\quad \beta'=p-1,
 \qquad
 C_1\longrightarrow C_0:\quad p'=\beta-1.
 \tag{R0E.11}
\]

For the first arrow keep original \(x=p^{-2}\) fixed and solve the
source equation analytically in \(t\). Its first derivative is

\[
 y_t(0)=-p/2,\quad x(p-G)=t\,p(p-2)+O(t^2),
 \quad y(p-G)=t(p-1)+O(t^2).
 \tag{R0E.12}
\]

The full second arrow has central

\[
 x(p-G)=(\beta-1)^{-2},\qquad
 y(p-G)=-(\beta-1)^{-3}.
 \tag{R0E.13}
\]

Substituting the full rational coefficient R0E.5 at both the original
and the translated source point gives

\[
 \left.\frac{B_{2,1}}t\right|_{C_0}
       =4\frac{p-2}{p},\qquad
 \left.\frac{B_{2,1}}t\right|_{C_1}
       =4\frac{\beta-1}{\beta-3}.
 \tag{R0E.14}
\]

The actual point \(2G\) has \(p=2\); \(O\) is \(p=0\).
The rational point section therefore has the following coefficients
in genuine nonzero line frames:

\[
 r_0=(p-2)/p,\qquad r_1=1,
 \qquad e_i=\sigma_2/r_i.
 \tag{R0E.15}
\]

For a raw coefficient \(B\) covering \(T\), the corresponding
nonzero-frame coefficient is exactly \(B\,r(Tp)/r(p)\).
Equations R0E.11,14,15 give \(4\) on both components. In particular
all four section points \(O,G,2G,3G\) remain in the module maps;
their apparent raw zeros or poles are their displayed generators.

\par

\medskip\noindent\textbf{R0E.3. The source graph that retains those modules}

\par\nobreak

At the root point \(O\), put \(\eta=-X/Y\) and

\[
 U=\eta^2X,\qquad
 U^2-U^3-\eta^4U-z^2\eta^6=0,\qquad U(0,0)=1.
 \tag{R0E.16}
\]

The two charts of the actual blowup \((\eta,z)\) are

\[
 \eta=zp:\quad\tau_{\rm old}=p;
 \qquad z=\eta v:\quad
 (x,y,t)=(v^2U,-v^3U,\eta^4v^4).
 \tag{R0E.17}
\]

The first point generator is exactly

\[
 \frac{p-2/(1+8z^4)}p.
 \tag{R0E.18}
\]

At \(T=(0,0)\) in the central good root, the equation gives

\[
 X(1+X^2)=Y^2-z^2.
 \tag{R0E.19}
\]

Thus the source graph is the blowup \((Y,z)\). On \(Y=z\beta\),

\[
 X=z^2\alpha,\qquad
 \beta^2=1+\alpha+z^4\alpha^3,
 \qquad (x,y,t)=(z^4\alpha,z^4\beta,z^4),
 \tag{R0E.20}
\]

which is exactly R0E.10. The other chart is \(z=Yv\),

\[
 X=Y^2C,\qquad C(1+Y^4C^2)=1-v^2,
 \qquad(x,y,t)=(Y^4v^2C,Y^4v^3,Y^4v^4).
 \tag{R0E.21}
\]

Every formula preserves the original \(t\) and the actual point
divisor. The source translation R0E.6 exchanges these two graph
centres and carries their strict point generators. Its inverse is
obtained by using \(+G'\). These are the graph and Cartier arrows
used for the next full finite-piece extension; the cap at original
infinity remains the complete K96 cap on its own actual domain.

\par

\medskip\noindent\textbf{R0E.4. The full M6 coefficient and the tangential source chart}

\par\nobreak

We now retain both actual data
\[
 (a,b;m)=(6,1;3l+2),\ l\geq0,
 \qquad (a,b;m)=(2,3;3l),\ l\geq1,
 \qquad n=\pm1,
 \tag{R0E.22}
\]
with every finite character \(\zeta_m^d\), the prescribed selected
fan and the whole K96 annular reference. In the second row the
suspension is the chronological third power of the one-step line
map. All root and source formulas in this section use the original
\(r=0\) source R0E.1.

The full Miller functions on that source are
\[
 \begin{gathered}
 s=1/4,\quad U=y-t-x/2,\quad V=y+t+x/2,\quad
 f_2=U/(x-s),\\
 f_3=\frac{s xU}{sV-(s+2t)x},\qquad
 f_6=f_3^2 f_{2,3G},\\
 x_3=8t+64t^2,\quad y_3=3t+96t^2+512t^3,\quad
 \mu_3=\frac{3x_3^2+t}{2y_3},\\
 x_6=\mu_3^2-2x_3,\quad
 y_6=\mu_3(x_3-x_6)-y_3,\qquad
 f_{2,3G}(p)=\frac{y-y_3-\mu_3(x-x_3)}{x-x_6}.
 \end{gathered}                                                   \tag{R0E.23}
\]
The quotients in R0E.23 are coefficients between their actual
Cartier modules. The numerator tangent and denominator vertical
line retain their point divisors also where individual displayed
rational functions have a zero or a pole. Direct cancellation in
\(\mu_3\) gives an analytic function at \(t=0\), with
\[
 \mu_3(0)=1/6,\qquad (x_6(0),y_6(0))=(1/36,-1/216),
 \qquad -x_6(0)/y_6(0)=6.
 \tag{R0E.24}
\]

Here is its entire quartic pullback. Put
\[
 N=zY-z^2-X/2,\quad u_3=8+64z^4,\quad
 v_3=3+96z^4+512z^8,
\]
\[
 \begin{split}
 \phi_2&=N/(z^2X-1/4),\\
 \phi_3&=\frac{XN}{zY+z^2-X/2-8z^4X},\\
 \phi_{23}&=
 \frac{zY-z^2v_3-\mu_3(z^4)(X-z^2u_3)}{z^2X-x_6(z^4)},
 \qquad \phi_6=\phi_3^2\phi_{23}.
 \end{split}                                                     \tag{R0E.25}
\]
Substitution in the original functions gives the exact identities
\[
 f_2=z^2\phi_2,\quad f_3=z^2\phi_3,\quad
 f_{2,3G}=z^2\phi_{23},\quad f_6=z^6\phi_6,
 \qquad
 \frac{B_{6,1}}{t^3}=\phi_2(p)^6\frac{\phi_6(Tp)}{\phi_6(p)}.
 \tag{R0E.26}
\]
In particular \(\phi_3|_0=X\), \(\phi_{23}|_0=6X\),
\(\phi_6|_0=6X^3\). The full source translation in R0E.6 therefore
gives central coefficient \(64\) in the actual Hom line. Each
\(\phi_2,\phi_3,\phi_{23},\phi_6\) changes sign under the quartic
deck. The coefficient in R0E.26 is invariant, including its source
argument. The powers \(z^2,z^6\) remain in the point-to-tensor
arrows throughout.

For the original two components the full functions give the following
limits. The first column uses R0E.9; the second uses R0E.10, with
\(\alpha\) the analytic solution of
\(\alpha+t\alpha^3=\beta^2-1\):
\[
\begin{array}{c|c|c}
 &C_0& C_1\\ \hline
 f_2&2/[p(p-2)]&t\,2(\beta-1)^2+O(t^2)\\
 f_3&p^{-2}&t(\beta-1)^3/(\beta-3)+O(t^2)\\
 f_{2,3G}&6/[p(p-6)]&6t(\beta-3)^2+O(t^2)\\
 f_6&6/[p^5(p-6)]&6t^3(\beta-1)^6+O(t^4).
\end{array}                                                       \tag{R0E.27}
\]
These are obtained from R0E.23 before multiplying the factors.
Combining them with the two directed translations R0E.11 gives
\[
 \left.\frac{B_{6,1}}{t^3}\right|_{C_0}
       =64\frac{p-6}{p},\qquad
 \left.\frac{B_{6,1}}{t^3}\right|_{C_1}
       =64\frac{\beta-1}{\beta-7}.
 \tag{R0E.28}
\]
For either even value \(a=2,6\), keep the actual point frames
\[
 r_{a,0}=(p-a)/p,\quad r_{a,1}=1,\quad
 e_{a,i}=\sigma_a/r_{a,i},\qquad
 \mathcal B_{a,1}=t^{-a/2}B_{a,1}\,r_a(Tp)/r_a(p).
 \tag{R0E.29}
\]
Equations R0E.14,28 give \(\mathcal B_{a,1}=2^a\) on both central
components. The chronological \(b=6/a\) product has central
coefficient \(64\). Its actual source map is \(p\mapsto p-bG\).
All four horizontal divisors \(O,aG,G,(a+1)G\), and all their
successive translates in the third power, remain in R0E.29.

The tangency point is also covered by an explicit source morphism.
On \(X\ne0\), put
\[
 T_s=t/x=z^2/X,\qquad Y_s=y/x=zY/X,\qquad
 x=Y_s^2-T_s-T_s^2=z^2X.
 \tag{R0E.30}
\]
The original resolution map is
\((x,Y_s,T_s)\mapsto(x,xY_s,xT_s)\), with the displayed smooth
source equation. Its two central components are \(T_s=0\) and
\(x=0\); their intersection has length two, since on \(T_s=0\)
the equation is \(x=Y_s^2\). R0E.30 includes the root points
\(X=\pm i,Y=z=0\). Its coordinates are invariant under the full
quartic deck. Together with R0E.17,20,21 these charts cover the
entire compact central root. At the tangency the coefficients in
R0E.14,28 tend respectively to \(4,64\), as \(p\) or \(\beta\)
tends to infinity. In the normalization used in EF.30b,
\(w=1/2-1/p\); its literal coefficient
\(128(3w-1)\) is exactly R0E.28, including this limit.

The root translation carries the graph itself. At \(O\), substitution
\(X=U/\eta^2,Y=-U/\eta^3\) in R0E.6 gives
\[
 X'=\eta(\eta-2z)/U+2z^2\eta^4/U^2,\qquad
 Y'=z+(1/\eta-z\eta^2/U)X'.
 \tag{R0E.31}
\]
The last expression is holomorphic and has first jet \(\eta-z\),
so \((Y',z)=(\eta,z)\) as analytic ideals. At \(T\) write
\[
 A=1-\frac{X(Y-z)^2}{(1+X^2)^2},\qquad
 \eta'=-X'/Y'
  =\frac{(Y-z)A(1+X^2)^2}{A(1+X^2)^3-z(Y-z)^3}.
 \tag{R0E.32}
\]
Its denominator is a unit and \((\eta',z)=(Y,z)\). The inverse
\(+G'\) has the inverse ideal squares. These equations lift both
source maps to the two projective graph charts. On the remaining
\(X\)-unit chart they agree with the actual rational translation
and extend by their inverse morphisms on the smooth minimal source.

For \(a=2,6\), the full marked-section equation in the \(O\)-graph is
\[
 \tau(aG')=z\tau_{a,\rm old}(t),\quad
 \frac{\eta-\tau(aG')}{\eta}
 =\frac{p-\tau_{a,\rm old}(z^4)}p
 =1-v\tau_{a,\rm old}(z^4),
 \tag{R0E.33}
\]
where the last expression is on the other projective chart,
\(\tau_{2,\rm old}=2/(1+8t)\) and
\(\tau_{6,\rm old}=-x_6/y_6\). Both are analytic, with limits
\(2,6\), and both marked sections lie in this \(O\)-chart.
The numerator and denominator in R0E.33 have equal exceptional
orders; their two strict directions remain distinct. The root point
line and the pullback of the original point line are thus identified
by these actual ratios on the common graph. On the \(T\)-graph and
R0E.30 neither marked section is present, so the coefficient is one.
All these maps are quartic-deck invariant.

The original III component matrix is
\(\left(\begin{smallmatrix}-2&2\\2&-2\end{smallmatrix}\right)\)
and its primitive kernel is \((1,1)\). The even marked point \(aG\)
and \(O\) lie on the identity component, so the line has degree zero
on each component. The Hom multidegrees are zero. R0E.14,28 give
order \(a/2\) on each component. After removing precisely that
whole-fibre factor, the rational Hom section has no vertical divisor;
its horizontal divisor is exactly canceled by the two modules in
R0E.29,33. On the smooth resolved source, a rational section with
zero divisor and its reciprocal extend in every local line frame.
This proves the full map through the tangent point and every marked
section, with no omitted codimension-two extension. In particular
\[
 U_G(p,v_M)=(p+G,B_{a,1}(p+G)^{-1}v_M),\qquad
                 t^{a/2}U_G
 \tag{R0E.34}
\]
is an actual invertible finite-source line map. The inverse always
uses the reciprocal at its actual inverse source point.

\par

\medskip\noindent\textbf{R0E.5. The entire period module and the intrinsic original base}

\par\nobreak

The completed objects in R0E.22 are the central fibres of the actual
K96 disc family. F46 gives their integral sphere topology on its
reference fan; the finite common blowups, both ruling maps and the
primitive integral Gysin and van Kampen maps of F35 carry that topology
to each prescribed lattice-basis fan. The root and annular maps used
in this passage are stated again in R0E.8 below, specialized to
\(r=0\). Thus the constructed compact threefold has
\(H^2(X;\mathbb Z)=H^4(X;\mathbb Z)=0\). Its scalar belongs to the
actual common contraction domain, and its full character and reference
remain fixed. Its ordinary fibre is the compact complex two-torus
obtained from the given point-line suspension.

A horizontal effective divisor would restrict to a nonzero effective
divisor on a generic two-torus. Its class is the restriction of a
class in \(H^2(X)=0\). For a flat Kähler form \(\omega\) on that
torus the divisor formula would give
\(0=\int c_1(D)\wedge\omega=\int_D\omega>0\), with all positive
multiplicities included. Consequently every polar divisor of a
meromorphic function is vertical. On a regular compact torus the
function is holomorphic after removal of codimension-two indeterminacy
and hence constant. It descends to the regular base. At a generic
smooth point of a reduced special component choose the exact
transverse coordinate \(t-t_0=u^e\), after taking the extending root
of its base unit. The finite Laurent principal part is invariant
under \(u\mapsto\zeta_eu\), since the function is constant on its
generic fibres. All exponents are divisible by \(e\), proving the
meromorphic extension on the base. This proves the actual pullback
\(\mathbb C(t)\xrightarrow{\pi^*}\mathbb C(X)\) is an isomorphism.
An endpoint isomorphism therefore has
\(\pi_{\rm target}F=h\pi_{\rm source}\) for a Möbius map \(h\),
with the same equality across all special fibres.

We identify the full source quotient before restricting this base map.
On a regular source disc choose its actual Abel representative
\(g\), period \(\tau\), and write
\[
 Z=e^{2\pi iz},\quad Q=e^{2\pi i\tau},\quad P=e^{2\pi ig},
 \quad f_P(Z;Q)=\frac{\Theta(Z/P;Q)}{\Theta(Z;Q)},\quad
 V_{\rm pt}=f_{P^a}(Z;Q)v_M.
 \tag{R0E.35}
\]
Here \(v_M\) is the coefficient in the canonical meromorphic point
section; its actual two Cartier modules, and the reciprocal arrow in
each, define R0E.35 through the section points. The theta product
satisfies \(f_P(QZ)=P f_P(Z)\). In particular the four periods of
the actual completed torus, in the order
\((\gamma,e_1,e_2,\delta)\), are
\[
 \gamma=(-bg,\ell),\quad e_1=(1,0),\quad
 e_2=(\tau,ag),\quad\delta=(0,1).
 \tag{R0E.36}
\]
The entire suspension multiplier defining \(\ell\) is
\[
 \lambda\frac q{t^3}B_{a,b}(p)
       \frac{f_{P^a}(P^{-b}Z;Q)}{f_{P^a}(Z;Q)}.
 \tag{R0E.37}
\]
Its ratio with the translation lift having the same constant factors
is an invertible function on the compact elliptic fibre, so it is a
base constant. R0E.37 specifies that whole constant, including every
Miller and selected-height factor. The tensor-frame dictionary is
\[
 V_{\rm tensor}=f_P(Z)^a f_a(p)^{-1}v_M
        =C_a(t)V_{\rm pt},\qquad
 C_a(t)=\frac{f_P(Z)^a}{f_a(p)f_{P^a}(Z)}.
 \tag{R0E.38}
\]
The full divisor identity proves independence of the source point;
its value and its logarithm continuation remain in \(\ell\).

For completeness the exact period-coordinate changes are recorded.
For \(g'=g+\mu+\nu\tau\), \(\mu,\nu\in\mathbb Z\), the entire
point coefficient changes by
\[
 V'_{\rm pt}=(-Z/P^a)^{a\nu}
              Q^{-a\nu(a\nu+1)/2}V_{\rm pt}.
 \tag{R0E.39}
\]
Its inverse pulls the new vectors to
\[
 e'_1=e_1-a\nu\delta,\quad e'_2=e_2+a\mu\delta,\quad
 \gamma'=\gamma-b\mu e_1-b\nu e_2+c\delta,\quad
 c=\ell'-\ell+ab\nu(2g+\mu+\nu\tau)\in\mathbb Z.
 \tag{R0E.40}
\]
For \(A=\left(\begin{smallmatrix}p&q_0\\r_0&s_0\end{smallmatrix}\right)
\in SL_2(\mathbb Z)\), set \(L=p+r_0\tau\). The full source/line map is
\[
 z'=z/L,\quad \tau'=(q_0+s_0\tau)/L,\quad g'=g/L,
 \quad w'=w-a r_0gz/L+c_A,
 \tag{R0E.41}
\]
where its entire constant is fixed by
\[
 e^{2\pi i c_A}=
 \frac{f_{e^{2\pi iag'}}(e^{2\pi iz'};e^{2\pi i\tau'})}
 {f_{P^a}(Z;Q)}\exp(2\pi i a r_0gz/L).
 \tag{R0E.42}
\]
The divisor and both period factors cancel, making this the stated
constant; evaluating it at an ordinary point fixes its full value.
The inverse solves R0E.41 with its actual target parameters. The
vectors \(p e_1+r_0e_2,q_0e_1+s_0e_2\) go exactly to the two new
elliptic periods. The resulting entire lattice monodromy is
\[
 M_\alpha=\begin{pmatrix}
 1&0&0\\ b v_\alpha&A_\alpha&0\\
 c_\alpha&a v_\alpha^t J A_\alpha&1
 \end{pmatrix},\quad J=\begin{pmatrix}0&-1\\1&0\end{pmatrix},
 \tag{R0E.43}
\]
\[
 v_{\alpha\beta}=v_\alpha+A_\alpha v_\beta,\qquad
 c_{\alpha\beta}=c_\alpha+c_\beta+
                  ab v_\alpha^tJA_\alpha v_\beta.
 \tag{R0E.44}
\]
The integral central cochain is the actual one from the full functions;
none of its entries is prescribed to vanish.

The rational Abel class in R0E.43 is nonzero. Indeed a rational
coboundary, multiplied by an integer and its full period cochain,
would give a single-valued holomorphic Lie logarithm of \(NG\).
Multiply \(N\) further by six to put it in every identity component.
At the III fibre R0E.4 puts \(2G'\) in the formal neighbourhood of
\(O\), with \(\tau(2G')=2z/(1+8z^4)\). The root formal logarithm
has deck weight \(i\); its quotient by \(z\), by R0E.3, is an
extending logarithm in the original finite Lie frame. The difference
from the proposed single-valued logarithm is an invariant elliptic
period. The central quartic action on \(dX/(2Y)\) has eigenvalue
\(i\), so it has no nonzero invariant period. At the selected I1,
the invariant annular period extends, by integration of the actual
invariant differential in \(xy=q\).

At original infinity the exact good cubic is
\(Y^2=X^3+zX+1\). The section \(3G'\) has root logarithm
\(-z^2/8+O(z^5)\). The original minimal differential is
\(\omega_0=z^{-2}\omega_{\rm good}\); thus division by \(z^2\)
gives its extending original logarithm. There is no invariant
elliptic period under its order-three central action. Consequently
the proposed logarithm is a global section of the actual Lie line.
Since \(x_0=q_\infty^2x,y_0=q_\infty^3y\),
\(\omega_0=q_\infty^{-1}\omega_t\). Its coefficient transition
is \(u_0(q)=q^{-1}u_t(q^{-1})\). An entire \(u_t\) satisfying this
is \(O(t^{-1})\) at infinity, hence zero. The Lie line is
\(\mathcal O(-1)\) by this explicit transition. On the other hand
\(3G\) has original central additive coordinate
\(-x_0/y_0=-1/8\); all its nonzero integral multiples are nonzero.
Thus \(G\) is not torsion and the rational Abel class cannot vanish.

The invariant integral subgroups of R0E.43 are now determined:
\[
 \Lambda^{\pi_1}=\mathbb Z\delta,\qquad
 (\Lambda^*)^{\pi_1}=\mathbb Z\psi,\quad
 \psi(\gamma)=1,\quad\psi(e_1)=\psi(e_2)=\psi(\delta)=0.
 \tag{R0E.45}
\]
An invariant vector with \(\gamma\)-coefficient \(z\) would solve
\((A_\alpha-I)u+bzv_\alpha=0\). The nonzero Abel class forces
\(z=0\), and infinity has no elliptic invariant vector, so \(u=0\).
The dual equation gives the second assertion. Primitivity follows
from the literal nonzero-line cover with lattice
\(\langle e_1,e_2,\delta\rangle\), whose primitive extra deck is
\(\gamma\).

An endpoint intertwines these invariant groups even before the base
map is fixed. On each regular fibre its lifted derivative is periodic
for the complete lattice, hence constant. Complex linearity and
\(F_*\delta=\epsilon\delta\) give the actual affine map and inverse
\[
 (z,w)\mapsto(Az+z_0,Cz+\epsilon w+w_0),\qquad
 (z',w')\mapsto
 ((z'-z_0)/A,\epsilon(w'-w_0-C(z'-z_0)/A)).
 \tag{R0E.46}
\]
Here \(A\ne0\) and \(\epsilon=\pm1\). The preserved subgroup
\(\ker\psi\) gives a degree-one elliptic source map, with complete
integer matrix
\[
 \begin{pmatrix}p&q_0&0\\ r_0&s_0&0\\u&v&\epsilon\end{pmatrix},
 \quad A=p+r_0\tau(h(t)),\quad
 A\tau(t)=q_0+s_0\tau(h(t)),\quad ps_0-q_0r_0=1.
 \tag{R0E.47}
\]
The determinant sign is positive because the first coordinate is a
nonzero complex-linear map. Thus the two period lattices are mapped
bijectively, with the displayed inverse. The original elliptic
quotient is identified through the actual point-line cyclic cover,
including all integers \(u,v,c_\alpha\).

Directly for R0E.1 the elliptic invariant is
\[
                         j_0(t)=\frac{6912}{27t+4}.
 \tag{R0E.48}
\]
R0E.46-\/-47 imply \(j_0(h(t))=j_0(t)\); substitution in this rational
function gives \(h(t)=t\). This treats every permitted base map,
including a proposed interchange of marked branches, using the actual
elliptic isomorphism just obtained.

\par

\medskip\noindent\textbf{R0E.6. Full linear centralizer and the primitive source section}

\par\nobreak

The rational elliptic monodromy centralizer is \(\mathbb QI\).
Its infinity matrix generates the quadratic field with polynomial
\(x^2+x+1\). The selected I1 has a nonzero unipotent nilpotent part,
so it cannot belong to that field. A non-scalar element of that field
generates it, proving the assertion for both monodromies together.
A matrix commuting with the full R0E.43 preserves R0E.45 and has form
\[
 X=\begin{pmatrix}s_0&0&0\\u&A&0\\c&w&t_0\end{pmatrix}.
\]
The middle block is \(A=eI\). Its middle-left equation is
\((I-A_\alpha)u=b(s_0-e)v_\alpha\), so the nonzero Abel class gives
\(s_0=e,u=0\). The lower-middle equation is the dual Abel coboundary
with coefficient \(a(e-t_0)\), giving \(t_0=e,w=0\). All remaining
central-cochain equations are satisfied exactly by
\[
                             X=eI+c\delta\psi.
 \tag{R0E.49}
\]
Integral unimodularity gives \(e=\pm1,c\in\mathbb Z\). Its complex
linear realization excludes \(c\ne0\): subtracting \(eI\) would
have real rank one, whereas a nonzero complex-linear map has even
real rank. Thus an endpoint has linear part precisely \(I\) or
\(-I\), with its complete affine displacement still retained.

For an automorphism of a single completed object the negative part
is excluded geometrically. It has sixteen transverse fixed points
on each regular two-torus, since its vertical fixed-point equation
is multiplication by two. The actual global equalizer with identity
is a closed analytic set in the compact threefold. An irreducible
component through a regular fixed point has dimension one and is the
proper analytic closure of a horizontal fixed curve. Its normalization
maps to the original base with positive degree. Pairing it with
\(\pi^*[\omega_{\mathbb P^1}]=0\) in \(H^2(X)\) gives zero, while
integration of this positive base form gives positive degree. This
contradiction uses the full cap-completed map. Negative isomorphisms
between two scalar cut endpoints will be treated by their actual
annular continuation in R0E.10.

An identity-linear endpoint lifts on the regular locus to the cyclic
nonzero-line cover, commuting with its actual generator \(B_\lambda\).
At the ordinary III disc, R0E.26-\/-34 give that cover as the invertible
contraction on the actual resolved point line. A nonzero line vector
over the \(O\)-section on the base disc kills the base meridian.
Generic-fibre loops therefore generate after the actual vanishing-cycle
relations, all in \(\ker\psi\). The preserved fibre homomorphism
extends and the covering-space lift is single valued through that
whole disc. Its source projection evaluated at this \(O\)-vector
is an extending section. The same argument applies to any ordinary
smooth or nodal source disc, retaining both Cartier modules.

At the selected node write its entire periods
\[
 h:(Z,V,q)\mapsto(QZ,\alpha V,q),\qquad
 B:(Z,V,q)\mapsto(gZ,q\eta V,q),\quad
 Q=\kappa q,\quad\alpha=P^a,\quad g=P^{-b}.
 \tag{R0E.50}
\]
All four displayed coefficient functions are the actual nonzero
units of R0E.35-\/-38. The source normalization of the selected fibre
is \(w=y/(x-2/9)\), with branches \(\pm\sqrt{2/3}\). Taking
\(Z=(w-\sqrt{2/3})/(w+\sqrt{2/3})\) gives
\[
             P(0)=\frac{2-\sqrt6}{2+\sqrt6}=-5+2\sqrt6.
 \tag{R0E.51}
\]
The full functions \(P(q),Q(q)\) remain in R0E.50. In a cone chart
\(q=z_0z_1z_2\), the image of a section has nonnegative coordinate
orders summing to one. Its two torus characters thus have integral
orders. The actual height vectors of \(h\) and \(B\) are the two
primitive coordinate generators, so a literal \(h^rB^k\) reindex
removes them. Keeping all its unit factors, the endpoint and inverse
then give extending units \(K(q),C(q)\). On every character their
map and inverse are
\[
 Z^uV^vq^d\mapsto K(q)^uC(q)^v Z^uV^vq^d,
 \qquad Z^uV^vq^d\mapsto K(q)^{-u}C(q)^{-v}Z^uV^vq^d.
 \tag{R0E.52}
\]
They are actual maps on every prescribed cone, with \(q\) unchanged.
The integer \(B^k\) is retained for the divisor equation below.

At infinity use the literal normalized fibre product with
\(z^m=q_\infty\). The maps in R0E.8 identify it with the root
quotient by \(B_\lambda\), since \(A_\lambda^m=B_\lambda^n\)
and \(n=\pm1\). Its canonical endpoint lift fixes \(z\).
The nonzero-line lift has commutator \(B_\lambda^k\) with the full
cap generator. Iterating \(m\) times gives \(B_\lambda^{mk}=1\)
because the lift already commutes with \(B_\lambda^n\). Strict
contraction gives infinite order, hence \(k=0\). This retains the
finite character. Evaluation extends the source section through
the actual root and its Cartier graph. The local lifts away from
the selected point differ by integer suspension powers. Their
orbifold base is \(\mathbb P^1\setminus\{t_c\}\) with its one
order-\(m\) cap. Its fundamental group is finite cyclic, so its
homomorphism to the integer deck group is zero. Removing this
cocycle by its actual deck powers gives a global source section.

We identify its section group integrally on this source. The rational
chart \(UV=x^2(x-1/4)\), with
\(t=(V-U-x)/2\), is the original surface. Resolving its A1 point
gives the rational elliptic surface. The actual fibres are IV*, III
and I1, with Euler sum \(8+3+1=12\); its intersection lattice is
unimodular of rank ten. In the A1 tangent cone
\(UV=-x^2/4\), the section \(G\) has direction
\([x:U:V]=[0:0:1]\), off the strict fibre plane \(V-U-x=0\).
It therefore meets the nonidentity III component once, away from
the tangency; \(O\) meets the identity component. The two sections
are disjoint. At infinity the actual cubic root has the distinct
fixed sections \((0,1)\) and \(O\), with quotient weights \((2,1)\)
and rays \((2/3,1/3),(1/3,2/3)\). Their two resolved arms are
disjoint. Label the E6 components as in the root charts, with edges
\((0,1),(1,2),(2,3),(3,4),(2,5)\) and \(G\) meeting component 0.
The positive Cartan matrix is
\[
 A=\begin{pmatrix}
 2&-1&0&0&0&0\\-1&2&-1&0&0&0\\
 0&-1&2&-1&0&-1\\0&0&-1&2&-1&0\\
 0&0&0&-1&2&0\\0&0&-1&0&0&2
 \end{pmatrix},\quad
 A^{-1}e_0=(4/3,5/3,2,4/3,2/3,1)^t.
 \tag{R0E.53}
\]
The nonidentity III block is \((2)\), with intersection vector 1.
The two corrections are \(4/3,1/2\); \(G\cdot O=0\), so
\(h(G)=2-4/3-1/2=1/6\). Both section self-intersections are \(-1\)
from the already identified Lie line. In the order
\((O,F,C_0,\ldots,C_5,C_{\rm III},G)\), take
\(O^2=G^2=-1,OF=GF=1\), block \(-A\oplus(-2)\),
\(GC_0=GC_{\rm III}=1\), with other unspecified pairings zero.
The determinant of this integral Gram matrix is \(-1\). The ten
actual divisor classes give the full lattice; restriction to the
generic fibre identifies its quotient by the first nine classes
with the free section group generated by \(G\). Equivalently the
height equations give \(\iota^2|\mathrm{MW}_{\rm tors}|^2=1\).
Thus the extending section above is exactly \(jG\), \(j\in\mathbb Z\).

\par

\medskip\noindent\textbf{R0E.7. The actual full automorphism factors through the physical scalar}

\par\nobreak

For \(a=2,6\), retain the entire point-line multiplier
\[
 B_{a,1}(p)=f_2(p)^a\frac{f_a(p-G)}{f_a(p)},\qquad
 B_{a,b}=U_G^{-b},\qquad B_{\rm sel}=(q/t^3)B_{a,b}.
 \tag{R0E.54}
\]
An ordinary endpoint lift covering \(+jG\) differs from the actual
\(U_G^j\) by a base function \(H(t)\), because their ratio on a
compact elliptic fibre is a holomorphic nonzero function. Its
extension is meromorphic in the actual finite and root Cartier
modules. R0E.52 gives its integral order at the selected node, so
\(H\) is a rational base function.

The exact coefficient orders, with the root order measured in
\(z^m=q_\infty\), are
\[
\begin{array}{c|ccc}
 &t=0&z=0&q=0\\ \hline
 U_G&-3/b&2m/b&0\\
 H&3j/b&-2mj/b&-k.
\end{array}                                                       \tag{R0E.55}
\]
The first column is R0E.26-\/-34, including the tangency and both
Cartier modules. The root column follows from the entire pullbacks
\[
 \begin{array}{c|c|c}
 (a,b;m)& B_{a,1,\rm old}&\text{central root Hom coefficient}\\ \hline
 (2,3;3l)&z^{-2l}B'_{2,1}&B'_{2,1}=-X\text{ in the canonical point frame}\\
 (6,1;3l+2)&z^{-6(l+1)}B'_6&B'_6/z^2\longrightarrow4 .
 \end{array}                                                       \tag{R0E.56}
\]
The root graph and strict-section arrows giving these expressions
are retained in R0E.8. Their generic orders do not replace those
extension conditions. The third column is the literal \(B^k\)
reindex in R0E.52. At every other finite irreducible fibre the Hom
numerator and denominator have equal generic order; none of the
remaining Miller factors contains that fibre. Thus there are no
other base divisors. Their degree equation is
\[
 3j/b-2j/b-k=0,\qquad j=bk,\qquad
 H=c(t^3/q)^k,\qquad H U_G^{bk}=c B_{\rm sel}^{-k}.
 \tag{R0E.57}
\]
The full quotient deck is \(B_\lambda=\lambda B_{\rm sel}\), so
this last map is exactly multiplication by \(c\lambda^k\) on
the actual point line. It extends through the original source,
root, fan and cap by their complete coefficient maps and commutes
with their annular square. Conversely this physical action gives
those actual holomorphic automorphisms. Its kernel is trivial:
an identity physical scalar would equal one fixed integral deck
power on the connected nonzero-line cover; that power translates
the source by \(-kbG\). The proven non-torsion of \(G\) forces
\(k=0\), and its scalar is one. Therefore on every object R0E.22,
\[
                  \operatorname{Aut}_{\rm hol}(X)
                        =\mathbb C^*_{\rm physical}.
 \tag{R0E.58}
\]
The group on the right is connected by its actual scalar paths.
All selected fan markings and finite character data in this equality
are the ones specified at the outset.

\par

\medskip\noindent\textbf{R0E.8. Both full infinity roots and the retained annular reference}

\par\nobreak

We give the explicit specialization of K96 rather than extending the
finite quartic root to infinity. In the reduced row put
\(m=3l+2,k=2l+1\). The full original-coordinate map is
\[
 \begin{gathered}
 q_\infty=z^m,\quad t=z^{-m},\quad
 x=z^{-2(l+1)}X,\quad y=z^{-3(l+1)}Y,\\
 Y^2=X^3+z^{l+2}X+z^2,\quad G'=(0,z).
 \end{gathered}\tag{R0E.59}
\]
Its deck is \(D(z,X,Y)=(\zeta_mz,\zeta_m^{-2k}X,
\zeta_m^{-3k}Y)\). With \(c=(1+z^lX)^{1/2}\), chosen with
constant term one, the exact coordinates
\(U=Y-zc,V=Y+zc\) have \(UV=X^3\). The equation
\(2zc=V-U\) has its analytic inverse for \(z\). The three regular
A2 charts are
\[
 X=\xi\eta,\quad U=\xi^{3-i}\eta^{2-i},\quad
 V=\xi^i\eta^{i+1},\qquad i=0,1,2.
 \tag{R0E.60}
\]
In the middle chart \(X=ab,U=a^2b,V=ab^2\), the entire \(+G'\)
map is
\[
 \begin{gathered}
 K=1+\frac{b(b-a)z^l}{2c(c+1)},\quad
 F=1-\frac{abz^l(K+1)}{2c(c+1)},\quad X'=a(a-b)F,\\
 c'=(1+z^lX')^{1/2},\quad
 H=K+\frac{b(a-b)z^l}{2c(1+c')},\\
 a'=(b-a)F/H,\quad b'=-aH.
 \end{gathered}                                                   \tag{R0E.61}
\]
It is the chord with slope \((Y-z)/X=aK\),
\(Y'=-z-aKX'\); hence \(V'=-aX'H\) and the inverse of
\(2zc=V-U\) retains exactly the same \(z\). The inverse chord
is obtained from \(Y\mapsto-Y\). At the other smooth group points
it is the ordinary translation, giving the inverse morphisms in
all three charts. The vertical field
\[
 W=a(2b-a-z^{l+1}/c)\partial_a+
   b(2a-b+z^{l+1}/c)\partial_b
 \tag{R0E.62}
\]
is holomorphic and vanishes at the triple point. Its flow is the
identity-component group flow on the dense smooth locus. Thus
\[
 R_{\rm div}=2G'-[m]^{-1}((2m-1)G'),\qquad mR_{\rm div}=G'
 \tag{R0E.63}
\]
is an actual extending translation: \(2m-1=3(2l+1)\) puts the
argument in the small identity formal group. The formal inverse
has derivative \(1/m\). The deck commutes with these arrows:
\(l-2k=-m\) makes \(c\) invariant; the two middle coordinates
have weight \(\zeta_m^{l+1}\), and every \(K,F,H\) is invariant.
Uniqueness of the formal inverse gives the same commutation for
R0E.63.

The point-line graph retains
\[
 \tau'=-X/Y=z^k\tau_0,\qquad
 \tau'(6G')=z^kL(q_\infty),\quad L(0)=-1/4.
 \tag{R0E.64}
\]
It is the normalized graph of \((\tau',z^k)\), resolved by the
ordered rays \((j,1)\), \(1\le j\le k\), together with its axes.
Adjacent determinants are one; both section equations have order
\(j\) on the ray \((j,1)\). On the equality ray \(j=k\) their
strict directions are separately \(0,L\). Their exact ratio is
\[
 \frac{\tau'}{\tau'-z^kL}=\frac u{u-L},\qquad
 \frac1{1-vL}\quad\text{on }\tau'v=z^k,\ uv=1.
 \tag{R0E.65}
\]
The remaining cyclic charts carry these same monomials and units.
The product of the actual commuting source/root translations
preserves the closure of the graph and its inverse. Consequently
\(\beta^*N_6=\pi^*M_6\) as actual Cartier maps. The full selected
coefficient is
\[
 B'_{\rm sel}=h z^{-2}B'_6,\quad
 h=1-t_cz^m=1+4z^m/27,\quad B'_6/z^2|_0=4.
 \tag{R0E.66}
\]
The last limit is the full K96 EX.13 chord calculation at
\(s=1/4\); all its factors are retained by R0E.23 and the graph.
Its zero vertical divisor after the displayed factor and its
literal horizontal Hom divisor give the extending map and inverse.

For the good row \(m=3l\), the original map is
\[
 q_\infty=z^{3l},\quad t=z^{-3l},\quad
 x=z^{-2l}X,\quad y=z^{-3l}Y,\qquad
 Y^2=X^3+z^lX+1,\quad G'=(0,1).
 \tag{R0E.67}
\]
Its deck is \((z,X,Y)\mapsto(\zeta_mz,\zeta_m^{-4l}X,Y)\).
Put \(u=z^l\). At \(2G'=(u^2/4,-1-u^3/8)\) use
\(w=X-u^2/4\); at \(O\) use \(w=-X/Y\). After the literal
\(\mu_l\) quotient these have weights \((2,1)\). The invariant
coordinates \((w^3,u^3,wu)\) have equation \(AB=C^3\), resolved
by the three charts of R0E.60 with these corresponding coordinates.
In the quotient lattice the two rays are
\((2/3,1/3),(1/3,2/3)\). The strict-section support is
\(\max(v_w-2v_u,0)\); its difference from the pulled-back root
section is \(\min(v_w,2v_u)\), hence
\[
 \beta^*N_2=\pi^*M_2+E_{2G'}-E_O,\qquad
 E_P(v_w,v_z)=\min(v_w,2l v_z).
 \tag{R0E.68}
\]
This equality uses the ratio of the actual monomial section
generators and their units on each cone, with inverse ratio in the
other direction. The small section
\(R_{\rm div}=[m]^{-1}(3G')\) has
\(\tau=-u^2/(8m)+O(u^5)\). Its translation has
\(w\mapsto w\varepsilon(w,u)+u^2a(w,u)\), with \(\varepsilon\)
a unit, and preserves \((w,u^j),(u,w^j)\), \(j=1,2\), with their
inverse graph arrows. Thus it carries both full corrections.
The original Miller factor is
\[
 f_{2,\rm old}=z^{-l}\phi,\quad
 \phi=\frac{Y-1-z^lX/2}{X-z^{2l}/4},\quad
 B'_{\rm sel}=hB'_{2,3},\quad B'_{2,3}|_0=4.
 \tag{R0E.69}
\]
Indeed on \(Y^2=X^3+1\) the three translated factors are
\((Y-1)/X,-(Y+1)/X,2/X\), whose product is \(-2\); their squared
norm is four. Their actual chord translations have order three
there. This verifies the entire central norm while preserving the
full functions away from the central curve.

Here are the exact norm lift and reference in both rows. For
\(N=\mathcal O(P_0-O)\), with \(P_0=aG'\), use the literal
chord-over-vertical coefficient
\[
 g_{P_0,R}(p)=\frac{\ell_{P_0,R}(p)}{x(p)-x(P_0+R)}.
 \tag{R0E.70}
\]
Its divisor is \((P_0)+(R)-(P_0+R)-(O)\), precisely the Hom
coefficient for the map covering \(p\mapsto p-R\). Set
\(F=\tau(R)g_{P_0,R}\) in the good row and
\(F=\tau(P_0)g_{P_0,R}\) in the reduced row. Their orders are
\(2l,k\), with leading coefficients \(-1/(8m),-1/4\).
Near \(O\), \(x=\tau^{-2}+O(1),y=-\tau^{-3}+O(\tau^{-1})\)
and the actual chord slope is \(-\tau^{-1}+O(1)\). Thus multiplying
that chord by the displayed \(\tau\) gives central coefficient
one. The horizontal Hom divisor and the primitive component kernel
extend \(F\) and its inverse through the entire root. Its chord
weight is \(\zeta_m^{-k}\), its \(\tau\)-factor weight
\(\zeta_m^k\), where \(k=2l\) in the good row. It therefore
commutes with the full deck on every Cartier chart.

The ratio \(\beta=B'_{\rm sel}/F^m\) covers identity and is an
invertible base function, central value four. It is deck invariant.
Its analytic \(m\)-th root is fixed by the full chosen reference;
its transform-to-itself ratio is a constant root of unity equal to
one at the central fibre. Define
\[
 \widetilde F=\beta^{1/m}F,\quad
 A_0=\zeta_m^d\widetilde F^{\,n}D,\quad
 R_\lambda^m=\lambda,\quad A_\lambda=R_\lambda^n A_0,\quad
 B_\lambda=\lambda B'_{\rm sel}.
 \tag{R0E.71}
\]
Then \(\widetilde F^m=B'_{\rm sel}\),
\(A_\lambda^m=B_\lambda^n\), and
\([A_\lambda,B_\lambda]=1\), as full source/line maps with
character \(\zeta_m^d\).

On the prescribed sectors set
\(s_\alpha=\log_\alpha q_\infty/(2\pi i)\),
\(z_\alpha=e^{2\pi i s_\alpha/m}\),
\(c_\lambda=n\log R_\lambda/(2\pi i)+d/m\), and
\(f_\alpha=\operatorname{Log}_\alpha\widetilde F\). The full map is
\[
 \Phi_\alpha=\operatorname{Exp}
 (n s_\alpha f_\alpha+s_\alpha c_\lambda\delta+h_\alpha)I_{z_\alpha}.
 \tag{R0E.72}
\]
It retains signed source displacement \(-n s_\alpha R_{\rm div}\),
scalar \(R_\lambda^{n s_\alpha}e^{2\pi i d s_\alpha/m}\), and the
literal root change. For \(s_\beta=s_\alpha+k\),
\(f_\beta=D_*^k f_\alpha+\Pi_{\alpha\beta}\), its complete
cochain equation is
\[
 h_\beta-D_*^kh_\alpha=-n s_\alpha\Pi_{\alpha\beta}
               +j_{\alpha\beta}\operatorname{Log}_\beta B_\lambda
               +\Lambda_{\alpha\beta}.
 \tag{R0E.73}
\]
All transported periods and all integers are those of the prescribed
K96 reference. In the fixed ordinary Lie frame take its extension
\(h^0\), residual \(E\), and fixed partition \(\chi_\beta\).
The explicit solve is
\(k_\alpha=-\sum\chi_\beta E_{\alpha\beta}\),
\(v(q)=\pi^{-1}\int\bar\partial k(w)/(q-w)\,dA(w)\),
\(h_\alpha=h_\alpha^0+k_\alpha-v\). The residual has the compact
support specified in K96 and vanishes at its reference; this solve
therefore preserves the full reference value and every integral
column. Subtracting the target-sheet logarithm
\(kn f_\beta+kc_\lambda\delta\) gives exactly
\[
 \Phi_\beta=A_\lambda^k B_\lambda^{j_{\alpha\beta}}\Phi_\alpha,
 \qquad \Phi_\alpha B_\lambda=B_\lambda\Phi_\alpha.
 \tag{R0E.74}
\]
The inverse is the inverse root map followed by the negative
translation at its actual inverse source point, with all reciprocal
line factors. These are the entire ordinary/cap gluing arrows used
above. Their finite roots extend by R0E.59-\/-71 on their own infinity
disc. The finite III quartic charts have original \(t\) near zero;
their domain comparison is precisely R0E.1 and R0E.59 or 67 on the
common ordinary annulus. Every base function and reference period
is pulled back by these stated original-coordinate maps.

\par

\medskip\noindent\textbf{R0E.9. The complete scalar family and its primitive positive endpoints}

\par\nobreak

Keep the full scalar root \(R^m=\lambda\) distinct from the source
coordinate \(t\). Starting at the prescribed reference, set
\[
 \lambda_\theta=e^{2\pi i\theta}\lambda_*,
 \quad R_\theta=e^{2\pi i\theta/m}R_*,
 \quad B_\theta=e^{2\pi i\theta}B_*,
 \quad A_\theta=e^{2\pi i n\theta/m}A_*.
 \tag{R0E.75}
\]
On its labelled reference sectors let \(J_\alpha\) be the actual
integer representative with \(J_\beta-J_\alpha=j_{\alpha\beta}\).
The entire scalar family is
\[
 \Phi_{\theta,\alpha}
 =e^{2\pi i n\theta s_\alpha/m}
    B_\theta^{J_\alpha}
      (B_*^{-J_\alpha}\Phi_{*,\alpha}).
 \tag{R0E.76}
\]
This formula retains the source-dependent periods \(\Pi,\Lambda\)
and their reference integers. On an overlap its two extra factors
are \(e^{2\pi i n\theta k/m}\) and
\(e^{2\pi i\theta j_{\alpha\beta}}\), giving exactly the
R0E.74 squares for \(A_\theta,B_\theta\). All reciprocal factors
give its inverse. At an integral winding \(J\),
\[
 B_{\theta+J}=B_\theta,\qquad
 \Phi_{\theta+J,\alpha}=z^{nJ}\Phi_{\theta,\alpha}.
 \tag{R0E.77}
\]
The retained suspension factors have ratio
\(e^{2\pi iJ J_\alpha}=1\), after their integers have been carried.

Every scalar endpoint fixes the original base by R0E.45-\/-48.
It intertwines the identical full monodromy matrices R0E.43,
since the scalar changes only the entire suspension logarithm.
Its lattice matrix has form R0E.49. The two complex-independent
vectors \(e_1=(1,0)\), \(\delta=(0,1)\) force its complex-linear
map to be \(eI\). Its suspension equation therefore says
\[
             e(\ell_{\rm target}-\ell_{\rm source})\in\mathbb Z.
 \tag{R0E.78}
\]
All source, height and Cartier factors in R0E.37 agree. Their
remaining difference is the logarithm of the original scalar
ratio, so the two scalar values are equal. Its scalar path has
integral winding \(J\), and its ordinary linear part is one of
the two signs. The integer in R0E.78 is retained as its suspension
lift change.

A positive endpoint has ordinary lift \(H(t)U_G^j\) by the complete
source/cover argument. R0E.77 gives its full extra root factor
\(z^{nJ}\). Consequently the three divisor equations become
\[
 \operatorname{ord}_0H=3j/b,\qquad
 \operatorname{ord}_{t_c}H=-k,\qquad
 m\operatorname{ord}_\infty H+2mj/b+nJ=0.
 \tag{R0E.79}
\]
There are no other base divisors, by R0E.26-\/-34 and 56.
Their degree yields
\[
                   j=bk+\frac{bnJ}{m},\qquad m\mid bJ.
 \tag{R0E.80}
\]
Thus the possible positive windings are multiples of \(m\) in the
M6 row and of \(l\) in the M2 row. We now construct their actual
primitive maps, including all cap characters.

For M2 the finite map \(S=tU_G\) extends by R0E.34. At the selected
node its full point coefficient is
\[
 C_S(t)=t\,B_{2,1}(p+G)^{-1}
              \frac{f_{P^2}(PZ;Q)}{f_{P^2}(Z;Q)}.
 \tag{R0E.81}
\]
It is constant on the compact source fibre by the full Hom argument,
and both it and its reciprocal extend at \(q=0\). Its actual
character map is \((Z,V)\mapsto(PZ,C_SV)\); on
\(Z^uV^vq^d\) the multiplier is \(P^uC_S^v\). It carries the same
fan with the unchanged height and complete units. Chronologically,
\[
                    S^3=t^3U_G^3=qB_{\rm sel}^{-1}.
 \tag{R0E.82}
\]
On the good root R0E.67, define \(U'_G=(B'_{2,1})^{-1}\) with
\(B'_{2,1}(p)=\phi(p)\phi(p-G')\) and the full \(\phi\) of
R0E.69. The coefficient of \(B'_{2,1}(p)\) at the central curve is
\(-X(p)\) in the canonical point frame. The actual inverse is
\[
 U'_G(p,v_M)=\bigl(p+G',B'_{2,1}(p+G')^{-1}v_M\bigr),
 \qquad B'_{2,1}(p+G')^{-1}\big|_0
   =-\frac1{X(p+G')}=\frac{Y+1}{2X}.
\]
Here \(X(p+G')=2(1-Y)/X^2=-2X/(Y+1)\); thus the coefficient
is evaluated at the literal inverse source point. The actual
Hom Cartier modules supply both maps\textquotesingle{} zeros and poles. Their
generic central order is zero, so each Hom map and its inverse
extend everywhere on the root. The complete
ordinary/root coefficient square is
\[
 f_{2,\rm old}=z^{-l}\phi,\quad
 B_{2,1,\rm old}=z^{-2l}B'_{2,1},\quad
                     U_{G,\rm old}=z^{2l}U'_G .
 \tag{R0E.83}
\]
It is read on the common graph R0E.68 with both actual correction
monomials, source and translated point modules retained.

The full character computation is
\[
 \phi(Dp)=\zeta_m^l\phi(p),\quad
 B'_{2,1}(Dp)=\zeta_m^{2l}B'_{2,1}(p),\quad
 U'_G D=\zeta_m^{-2l}D U'_G.
 \tag{R0E.84}
\]
The source translations of \(U'_G\) and the full norm lift
\(\widetilde F\) commute. In a common constant-factor point frame
both line coefficients are base constants, specified by their
entire rational maps; their lifted affine translations thus
commute, including their complete values. The root \(\beta^{1/m}\)
is also a base function. With
\(A_R=R^n\zeta_m^d\widetilde F^{\,n}D\), this gives the exact square
\[
       (U'_G)^n A_R=\zeta_m^{-2ln}A_R(U'_G)^n
                    =\zeta_m^{ln}A_R(U'_G)^n .
 \tag{R0E.85}
\]

For the positive scalar winding \(J=l\), so \(R\mapsto\zeta_3R\),
the forward endpoint is
\[
 E_R^{\rm ord}=(R^{-l}tU_G)^n,\qquad
 E_R^{\rm cap}=(R^{-l}U'_G)^n.
 \tag{R0E.86}
\]
For \(n=-1\) these are the actual inverse maps with the reciprocal
at the inverse source point. They satisfy, with the entire
character \(d\),
\[
 E_R^{\rm cap}A_R=A_{\zeta_3R}E_R^{\rm cap},\qquad
 \Phi_{\theta+l}E_R^{\rm ord}\Phi_\theta^{-1}
 =z^{nl}R^{-ln}t^n z^{2ln}(U'_G)^n
 =E_R^{\rm cap}.
 \tag{R0E.87}
\]
The last cancellation uses precisely the original \(t=z^{-3l}\).
Its inverse square retains the same units and all reference
periods. Therefore R0E.86 is a holomorphic endpoint on each
prescribed fan and on the full completed object.

The three chronological endpoints use
\(R,\zeta_3R,\zeta_3^2R\). Their scalar factors multiply to
\(\lambda^{-n}\), including the successive phases. Their full
composites are
\[
 \begin{split}
 E_{\zeta_3^2R}^{\rm ord}E_{\zeta_3R}^{\rm ord}E_R^{\rm ord}
     &=\lambda^{-n}S^{3n}=q^n B_\lambda^{-n},\\
 E_{\zeta_3^2R}^{\rm cap}E_{\zeta_3R}^{\rm cap}E_R^{\rm cap}
     &=\lambda^{-n}(U'_G)^{3n}=h^n B_\lambda^{-n}.
 \end{split}                                                     \tag{R0E.88}
\]
Thus on the actual quotient their endpoint is the canonical
\(q^n/h^n\) pair. In the M6 row this pair itself is the primitive
endpoint at \(J=m\): R0E.77 gives
\(z^{nm}q^n=(q_\infty q)^n=h^n\) on the overlap, with inverse
at the return. Its selected zero is carried by the actual
\(B_\lambda^n\) reindex and reciprocal full source/line unit.
Specifically in R0E.50,
\(q^n=B_\lambda^n(P^{bn},\eta^{-n})\); this latter ordered pair
means its actual source and line torus-unit map. Its character
arrows and inverses extend on every selected cone. Thus the zero
is retained by that literal quotient square, for either sign.
Chronological powers and inverses of these primitive maps give
every positive-part winding in R0E.80.

\par

\medskip\noindent\textbf{R0E.10. The full negative endpoint obstruction and exact winding group}

\par\nobreak

The negative argument uses the actual selected inversion and the
whole annular reference. On R0E.50 let
\[
                    I(Z,V,q)=(Z^{-1},V^{-1},q).
 \tag{R0E.89}
\]
It satisfies \(Ih=h^{-1}I\) and \(IB=B^{-1}I\), with all coefficient
functions retained at the unchanged \(q\). For a periodic fan of
basis \((u,v)\), negation sends its triangle
\(\operatorname{conv}(p,p+u,p+u+v)\) to
\(\operatorname{conv}(p',p'+v,p'+u+v)\), where
\(p'=-p-u-v\). This is the other triangle of the same periodic fan.
The full chart arrows can be read directly. With
\(\mathcal A=\kappa g,\mathcal C=\alpha\eta\), the representative
\(hBI\) sends
\[
 (a,b,c)\mapsto(\mathcal A\mathcal C a,\mathcal C^{-1}c,
                         \mathcal A^{-1}b),\qquad
 (u_0,v_0,w_0)\mapsto((\mathcal A\mathcal C)^{-1}u_0,
                          \mathcal A w_0,\mathcal C v_0).
 \tag{R0E.90}
\]
The first source chart has \(q=abc,Z=b,V=c\), the second
\(q=u_0v_0w_0,Z=u_0v_0,V=u_0w_0\). The displayed maps interchange
these charts and are literal inverses. Both preserve the full
product \(q\) and give \(Z'=\mathcal A q/Z,V'=\mathcal C q/V\).
For a general basis the actual integral character change carries
these equations with its complete period products. This proves
the inversion extends to the full selected quotient.

A negative endpoint lifts locally to the nonzero-line cover and
satisfies
\[
                        \widetilde F B_\lambda
                              =B_\lambda^{-1}\widetilde F.
 \tag{R0E.91}
\]
Any two local lifts differ by an integral suspension deck. At the
ordinary III disc the actual cover and Cartier maps are R0E.26-\/-34.
A nonzero line vector over the \(O\)-section of the entire disc
kills the base meridian. Generic fibre loops then generate modulo
the actual vanishing cycles in \(\ker\psi\). The preserved
negative cover class gives a single-valued holomorphic lift,
including its local biholomorphic inverse through the two
Cartier modules. At the selected point compose with R0E.89.
The result has positive linear part. R0E.50-\/-52 remove its
integer valuations using the entire \(h^rB^k\), and give a
single-valued bounded unit map and inverse on the punctured disc.
Composing back with R0E.89 gives a single-valued negative lift
there. The finite base with the selected point removed has
fundamental group generated by this one meridian. The ordinary
discs add no puncture. Continuing the local lifts therefore
gives one single-valued map on the whole ordinary nonzero-line
cover, retaining their constant joining deck integers.

On the actual normalized root fibre product, the kernel of the
finite generator is \(\langle B_\lambda\rangle\) because
\(A_\lambda^m=B_\lambda^n\), \(n=\pm1\). Its root nonzero-line
cover has the analogous \(O\)-section vector over the whole
\(z\)-disc. On a point/O intersection the frame is exactly
\(\sigma_a/[(\tau-\tau(aG'))/\tau]\), and elsewhere it is the
single-divisor frame; the transition is the actual Cartier unit
of R0E.65 or 68. The same meridian and vanishing-cycle argument
therefore supplies a single-valued root lift
\(\widetilde F_{\rm cap}\), fixing \(z\) and satisfying R0E.91.

Follow one positive original \(q_\infty\)-circle through every
reference sector in R0E.74. Its actual arrow is
\[
 A_\lambda B_\lambda^\nu,\qquad \nu\in\mathbb Z,\qquad
 (A_\lambda B_\lambda^\nu)^m=B_\lambda^N,\quad N=n+m\nu.
 \tag{R0E.92}
\]
The integer \(\nu\) is the retained sum of reference suspension
integers; all analytic source/line factors and period columns
are already in these full operators. No value is imposed on
\(\nu\). The scalar family R0E.76 retains it at both endpoints.
At winding \(J\), its extra \(z^{nJ}\) is single valued on the
root circle, so both endpoint annular maps have this same
monodromy.

Choose one constant suspension deck so that on a cut annulus
the actual lifts satisfy
\[
 \widetilde F_{\rm cap}
   =\Phi_{\rm target}\widetilde F_{\rm ord}\Phi_{\rm source}^{-1}.
 \tag{R0E.93}
\]
After one positive root circle its right side continues to
\[
 B_\lambda^N\widetilde F_{\rm cap}B_\lambda^{-N}
                   =B_\lambda^{2N}\widetilde F_{\rm cap}.
 \tag{R0E.94}
\]
Both actual lifts were proved single valued, so
\(B_\lambda^{2N}=1\). Strict contraction implies \(N=0\), whereas
\(n+m\nu=0\) has no integral solution for \(n=\pm1,m\ge2\).
This excludes every negative endpoint, including the quadratic
reduced cap, without dropping the reference integer.

The completed classification on the fixed original \(r=0\) source
is consequently
\[
 \begin{array}{c|c|c|c}
 (a,b;m)&\text{winding group}&\text{primitive ordinary forward map}
                                &\text{primitive cap map}\\ \hline
 (6,1;3l+2),\ l\ge0&m\mathbb Z&q^n&h^n\\
 (2,3;3l),\ l\ge1&l\mathbb Z&(R^{-l}tU_G)^n&(R^{-l}U'_G)^n .
 \end{array}                                                       \tag{R0E.95}
\]
It holds for every prescribed lattice-basis selected fan, both
norm signs and every retained finite character and reference.
The ratio of two endpoints with the same winding is exactly
the physical scalar group R0E.58. Each inverse return uses the
inverse of the displayed forward endpoint. R0E.88 is their
actual chronological threefold comparison. This proof assigns
no new numerical return.

\par

\medskip\noindent\textbf{R0E.11. Exact verification and preserved stages}

\par\nobreak

The original 47-check finite-chart stage, before this extension,
is preserved byte for byte under
\path{records/attempts/e6_iii_endpoint_before_global_extension/}.
Its recorded scope was R0E.1-\/-3. The added proof retains that
source map, proves the complete second marked-line and tangential
chart calculation, and supplies the entire global endpoint
classification in R0E.4-\/-10. The companion exact program checks
the full rational equations, ideal/chart and point-frame
identities, period-coordinate matrices, integral source lattice,
both root powers and characters, and the signed endpoint and
reference products. Its finite algebra is accompanied by the
explicit compactness, module-extension, covering and
single-valuedness arguments above. The final source and executed
certificate are submitted for independent full audit.

\par\medskip\noindent\textbf{Exact source and independent verification.}\par
The full mathematical source is
\path{notes/part3_e6_iii_fixed_source_endpoint_working.md}.
The primary exact certificate passes 9532 checks, including 116 full
integral elliptic basis changes and 4838 signed whole-annulus
instances. The independent finite-source proof and certificate pass
49 checks on the exact R0E.1--4 prefix:
\path{records/checks_audit_e6_iii_finite_source.json}.
All full-scope review gates and their final source locks are recorded
in \path{records/k114_export_manifest.json}.
The historical finite-stage files and the preserved algebraic-number
checker evaluation are retained in the two attempt directories named
in R0E.11. The original source bytes are unchanged by this export.

% End expanded verified cp36 input

\par\bigskip\noindent\textbf{Full independent root proof of the original III endpoint classification}\par\nobreak

% Expanded verified cp36 input: notes/part3_e6_iii_endpoint_root_peer.tex
We retain the original surface, line, selected height, finite character,
reference and fan of R0E.1-\/-95. In particular
\[
 y^2=x^3+tx+t^2,\quad G=(0,t),\quad q=t+4/27,
 \qquad (a,b;m)=(6,1;3l+2),\ l\ge0,
 \quad (2,3;3l),\ l\ge1,
 \tag{R0R.1}
\]
and the norm is \(n=\pm1\). The finite III disc retains the actual
order-one cap and zero norm of K96. Each scalar is in the actual contraction
domain of the complete K96 construction. The assertions below concern
these completed objects and their stated scalar paths. We prove their
full automorphism group and their scalar endpoint winding groups, with
the original-coordinate maps throughout. We assign no return integer.

\par

\medskip\noindent\textbf{R0R.1. Reconstructing the finite source and every point-line factor}

\par\nobreak

Put \(u=y-t-x/2,v=y+t+x/2\). Direct multiplication gives
\(uv=x^2(x-1/4)\), with inverse
\(y=(u+v)/2,t=(v-u-x)/2\). On \(x\ne0\), write
\(u=xu_1,v=xv_1\); then
\[
 x=u_1v_1+1/4,\qquad y=x(u_1+v_1)/2,
 \qquad t=x(v_1-u_1-1)/2.                       \tag{R0R.2}
\]
The reciprocals are the displayed \(u/x,v/x\). These are actual
birational source morphisms on their indicated open sets.

The A1 point of the original surface is resolved in its \(x\)-chart by
\[
 (x,Y_s,T_s)\longmapsto(x,xY_s,xT_s),\qquad
 x=Y_s^2-T_s-T_s^2.                              \tag{R0R.3}
\]
On \(x\ne0\) the inverse is \((x,y/x,t/x)\). The original fibre
function is \(t=xT_s\). Its two components have equations \(x=0\)
and \(T_s=0\); their intersection ideal has quotient
\(\mathbb C\{Y_s\}/(Y_s^2)\). Their intersection multiplicity is
two. Pairing the whole fibre with either component therefore gives
their intersection matrix \(\left(\begin{smallmatrix}-2&2\\2&-2
\end{smallmatrix}\right)\), with primitive kernel \((1,1)\).

The actual quartic source cover is
\[
 t=z^4,\quad x=z^2X,\quad y=z^3Y,\qquad
 Y^2=X^3+X+z^2,
 \quad D(z,X,Y)=(iz,-X,iY).                       \tag{R0R.4}
\]
The equation pulls back with the exact factor \(z^6\). The three
original coordinates are invariant under \(D\). In R0R.3 its map is
\(T_s=z^2/X,Y_s=zY/X\), defined on \(X\ne0\), including the two
central root points \((X,Y)=(i,0),(-i,0)\). Substitution gives
the full R0R.3 equation. Thus this chart includes the tangency in the
original fibre, with its full nonreduced intersection scheme.

We reconstruct the Miller functions from addition. The tangent slope
at \(G\) is \(1/2\). Its double and its chord with \(G\) give
\[
 2G=(1/4,-t-1/8),\quad L_{2G,G}=-8t-1/2,
 \quad 3G=(8t+64t^2,3t+96t^2+512t^3).             \tag{R0R.5}
\]
Write this last point as \((x_3,y_3)\), put
\(L_3=(3x_3^2+t)/(2y_3)\), and set
\(x_6=L_3^2-2x_3,y_6=L_3(x_3-x_6)-y_3\). The entire factors are
\[
 f_2=\frac{y-t-x/2}{x-1/4},\quad
 g_{2G,G}=\frac{y+t+1/8+(8t+1/2)(x-1/4)}{x-x_3},
 \quad f_3=f_2g_{2G,G},
 \tag{R0R.6}
\]
\[
 g_{3G,3G}=\frac{y-y_3-L_3(x-x_3)}{x-x_6},
 \qquad f_6=f_3^2g_{3G,3G}.                       \tag{R0R.7}
\]
Reduction by the original cubic identifies R0R.6 with the entire
\(f_3\) of R0E.23. Its divisor is \(3(G)-(3G)-2(O)\); the divisor
of R0R.7 is \(6(G)-(6G)-5(O)\). Thus these coefficients define
the required point-to-tensor Hom maps, with their actual point and
zero Cartier modules. The computation retains the rational functions
before taking any special-fibre value.

In R0R.4 the translation by \(-G'\), with \(G'=(0,z)\), is
\[
 L=(Y+z)/X,\quad X_T=L^2-X,
 \quad Y_T=z-LX_T.                               \tag{R0R.8}
\]
At the image point the slope for \(+G'\) is exactly \(-L\);
its two coordinates return \((X,Y)\). Both maps preserve the full
root equation. At \(O\), put \(X=U/\eta^2,Y=-U/\eta^3\), where
\(U^2-U^3-\eta^4U-z^2\eta^6=0,U(0,0)=1\). Then
\[
 X_T=\eta(\eta-2z)/U+2z^2\eta^4/U^2,
 \quad Y_T=z+(1/\eta-z\eta^2/U)X_T.              \tag{R0R.9}
\]
This last expression is holomorphic with first jet \(\eta-z\).
It identifies the ideals \((Y_T,z)=(\eta,z)\). At \(T=(0,0)\),
the expression R0E.32 has unit denominator and first jet \(Y-z\),
and identifies \((\eta_T,z)=(Y,z)\). Their inverse arrows have
the inverse ideal equalities. Hence these maps lift to both actual
blowups, in both directions, through every point of their exceptional
curves. R0E.17,20,21 and R0R.3 give the original-coordinate maps of
these charts. They cover the entire compact central root.

For \(a=2,6\), the marked point has
\(\eta(aG')=z\tau_{a,\mathrm{old}}(z^4)\), with finite limits
\(\tau_{2,\mathrm{old}}(0)=2\) and
\(\tau_{6,\mathrm{old}}(0)=6\). The two graph charts therefore
carry the exact point ratios
\[
 \frac{p-\tau_{a,\mathrm{old}}(z^4)}p,
 \qquad 1-v\tau_{a,\mathrm{old}}(z^4),
 \quad \eta=zp,\quad z=\eta v.                  \tag{R0R.10}
\]
Their inverses are the reciprocal Cartier-module arrows. Both marked
divisors have the same exceptional order; the two strict directions
remain separate. This proves the comparison of the two point lines
on the graph. The other graph centre and the tangency chart carry no
marked divisor, so their coefficient is one.

On the identity component use \((x,y)=(p^{-2},-p^{-3})\); on the
other use \((x,y)=(t\alpha,t\beta)\), with
\(\alpha+t\alpha^3=\beta^2-1\). R0R.8 gives the directed component
maps \(\beta'=p-1\) and \(p'=\beta-1\). Substituting the full
functions R0R.6-\/-7 first gives
\[
 \begin{array}{c|cc}
 &C_0&C_1\\\hline
 f_2&2/[p(p-2)]&2t(\beta-1)^2+O(t^2)\\
 f_3&p^{-2}&t(\beta-1)^3/(\beta-3)+O(t^2)\\
 g_{3G,3G}&6/[p(p-6)]&6t(\beta-3)^2+O(t^2)\\
 f_6&6/[p^5(p-6)]&6t^3(\beta-1)^6+O(t^4).
 \end{array}                                                       \tag{R0R.11}
\]
Consequently \(B_{a,1}=f_2^a f_a(Tp)/f_a(p)\), after the exact
whole-fibre factor \(t^{a/2}\), has central coefficients
\[
 2^a\frac{p-a}p\quad\hbox{and}\quad
 2^a\frac{\beta-1}{\beta-a-1}.                   \tag{R0R.12}
\]
Using the actual point-frame factors \(r_{a,0}=(p-a)/p,r_{a,1}=1\),
the Hom coefficient is the raw factor times \(r_a(Tp)/r_a(p)\).
It is \(2^a\) on both components. The exact graph generators in
R0R.10 retain this cancellation at the section points. The primitive
component kernel and the zero horizontal Hom divisor show that the
remaining rational section and its reciprocal have zero divisor on
the smooth resolved source. In each local regular ring they are
units, including at the tangency. This proves the complete finite
extension. Its chronological \(b=6/a\) product has coefficient 64.
The coefficient \(128(3w-1)\) in the earlier source chart equals
\(64(p-6)/p\) under the actual map \(w=1/2-1/p\).

\par

\medskip\noindent\textbf{R0R.2. The full invariant quotient, original base and primitive section}

\par\nobreak

The actual K96 disc family supplies the completed topology at this
source, with the full finite charts just proved and its own infinity
root, not the finite quartic root. The integral F35/F46 maps carry
\(H^2=H^4=0\) through each prescribed fan: their primitive normal
circle is unchanged by the common blowups and the two ruling maps.
On a compact scalar interval the original finite, root and annular
arrows have a common strict contraction; their compact gluing gives
the corresponding smooth transport. Thus these are the same sphere
objects used throughout R0R.1.

A horizontal effective divisor has positive pairing with a flat
Kähler form on a generic compact two-torus, but its class is the
restriction of a class in \(H^2(X)=0\). Hence all polar divisors of
meromorphic functions are vertical. Restriction to each regular
compact torus and removal of indeterminacy give constancy there.
At a special component the actual transverse map is
\(u\mapsto t_0+u^e\). Invariance of its Laurent expansion makes all
exponents multiples of \(e\), proving meromorphic descent on the
whole base. Thus \(\pi^*\mathbb C(t)=\mathbb C(X)\), and a
holomorphic endpoint isomorphism and its inverse induce a Möbius
base isomorphism \(h\).

We retain the complete point-theta module R0E.35-\/-44, whose periods
are \(\gamma=(-bg,\ell),e_1=(1,0),e_2=(\tau,ag),\delta=(0,1)\).
Its full monodromy is
\[
 M_\alpha=\begin{pmatrix}1&0&0\\bv_\alpha&A_\alpha&0\\
 c_\alpha&av_\alpha^tJA_\alpha&1\end{pmatrix},
 \qquad J=\begin{pmatrix}0&-1\\1&0\end{pmatrix}.   \tag{R0R.13}
\]
Both point Cartier modules, the full Miller-to-theta constant, and
every central cochain entry remain in this formula. Multiplication
gives \(v_{\alpha\beta}=v_\alpha+A_\alpha v_\beta\) and
\(c_{\alpha\beta}=c_\alpha+c_\beta+
abv_\alpha^tJA_\alpha v_\beta\). The companion checks these as
symbolic identities for arbitrary matrices of determinant one.

The Abel class is nonzero: a rational coboundary, after multiplication
by an integer and the finite component orders, gives a global
holomorphic Lie logarithm of a multiple of \(G\). At III the
logarithm of \(2G'\) divided by \(z\) extends by R0R.4-\/-5; the
quartic central action has no invariant elliptic period. At the
selected node its invariant annular period extends in \(xy=q\).
At infinity the actual good cubic is \(Y^2=X^3+zX+1\), the logarithm
of \(3G'\) has leading term \(-z^2/8\), and division by \(z^2\)
is the original Lie coefficient. Its order-three elliptic action
also has no invariant period. The full Lie transition is
\(u_0(q_\infty)=q_\infty^{-1}u_t(q_\infty^{-1})\). An entire
\(u_t\) satisfying it tends to zero at infinity and is zero.
But the original infinity value of \(3G\) is \((64,512)\), with
nonzero additive parameter \(-1/8\); every nonzero multiple remains
nonzero. This contradiction proves both non-torsion and the claimed
nonzero Abel class, with its actual finite extensions.

The vector and dual equations in R0R.13 now give exactly
\(\Lambda^{\pi_1}=\mathbb Z\delta\) and
\((\Lambda^*)^{\pi_1}=\mathbb Z\psi\), where
\(\psi(\gamma)=1\) and \(\psi(e_1)=\psi(e_2)=\psi(\delta)=0\).
Indeed any nonzero first vector coefficient makes the Abel cocycle a
coboundary; the remaining elliptic vector vanishes by the infinity
matrix. The dual equation has the same Abel obstruction. The extra
deck of the actual point-line cover is the primitive \(\gamma\).

An endpoint preserves these two integral subgroups with their signs
before fixing the base. The lifted derivative is periodic for the
whole compact torus lattice, so it is constant. Complex linearity
and \(\delta\mapsto\epsilon\delta\) give the complete affine map
\[
 (z,w)\mapsto(Az+z_0,Cz+\epsilon w+w_0),
 \quad (z',w')\mapsto
 ((z'-z_0)/A,\epsilon(w'-w_0-C(z'-z_0)/A)),         \tag{R0R.14}
\]
with \(A\ne0,\epsilon=\pm1\). The kernel cover is precisely the
actual nonzero point line, using both Cartier modules. On its lattice
the complete matrix is
\[
 K=\begin{pmatrix}p&q_0&0\\r&s&0\\u&v&\epsilon\end{pmatrix},
 \quad K^{-1}=\begin{pmatrix}s&-q_0&0\\-r&p&0\\
 \epsilon(vr-us)&\epsilon(uq_0-vp)&\epsilon\end{pmatrix},
 \quad ps-q_0r=1.                               \tag{R0R.15}
\]
The last determinant sign follows from the positive oriented real
area of a nonzero complex-linear map. The first coordinates of
R0R.14 and its inverse are actual inverse elliptic isomorphisms;
the integers \(u,v\) remain in the line arrow. For this exact
source \(j_0(t)=6912/(27t+4)\). Therefore
\(j_0(h(t))=j_0(t)\) gives \(h(t)=t\) on the regular locus and
hence on the whole base.

The full linear centralizer is \(eI+c\delta\psi\). For its elliptic
block, the order-three infinity matrix generates \(\mathbb Q(\zeta_3)\),
and the nontrivial selected unipotent matrix has a nonzero nilpotent
part. No nonzero nilpotent lies in that field. Every nonscalar field
element generates it, so the simultaneous elliptic centralizer is
\(\mathbb QI\). The two off-diagonal equations then use the nonzero
Abel class to set both adjoining scalars equal to \(e\) and both
elliptic vectors to zero. All central cochains remain arbitrary in
the surviving formula. Unimodularity gives \(e=\pm1\). Complex
linearity excludes the rank-one remainder when the fibre periods
coincide. A negative automorphism of one object has sixteen transverse
fixed points on each regular fibre. Its closed analytic equalizer in
the compact total space has a horizontal curve of positive base
degree. Pairing that curve with \(\pi^*H^2(\mathbb P^1)=0\) contradicts
that degree. Thus the automorphism has positive linear part.

The actual cyclic-cover, vanishing-cycle and selected-vertex unit maps
R0E.50-\/-52 extend its source displacement to a global source section.
At the finite III disc the nonzero \(O\)-section vector kills the
meridian, and all remaining vanishing cycles lie in \(\ker\psi\).
At the cap the normalized fibre product \(z^m=q_\infty\) has its
literal root lift. Its commutator with the cap generator is \(B^k\);
the \(m\)-fold relation makes \(B^{mk}=1\), so strict contraction
sets \(k=0\). The integer joining torsor on the one-punctured base
with one finite root has zero homomorphism to \(\mathbb Z\), since
its orbifold fundamental group is finite. These are the actual
single-valued lifts used to extend the section.

Its primitive generator is \(G\). The rational source R0R.2 has
fibres IV*, III and I1 and Euler number 12. The two sections are
disjoint, including their different infinity quotient arms. The
III A1 direction of \(G\) is \([x:u:v]=[0:0:1]\), off the strict
fibre plane \(v-u-x=0\), so its correction is \(1/2\). Its E6
arm correction is \(4/3\), giving height \(1/6\). More fully, the
ten actual classes \((O,F,C_0,\ldots,C_5,C_{\rm III},G)\), with the
pairings stated in R0E.53, have determinant \(-1\). They therefore
span the whole rank-ten unimodular divisor lattice. Quotienting by
the first nine classes gives the full free section group \(\mathbb ZG\),
with no torsion or index omitted. The extending section is \(jG\).

\par

\medskip\noindent\textbf{R0R.3. The actual line map and every positive scalar endpoint}

\par\nobreak

The lift over \(+jG\) is \(H(t)U_G^j\), where
\(U_G(p,v)=(p+G,B_{a,1}(p+G)^{-1}v)\). The full finite source
calculation gives its order \(-3/b\) at each III component. The
original infinity root maps are exactly
\[
 \begin{array}{c|c|c}
 m&t,x,y&\text{root equation}\\\hline
 3l+2&z^{-m},z^{-2(l+1)}X,z^{-3(l+1)}Y
     &Y^2=X^3+z^{l+2}X+z^2\\
 3l&z^{-3l},z^{-2l}X,z^{-3l}Y
     &Y^2=X^3+z^lX+1.
 \end{array}                                                       \tag{R0R.16}
\]
The reduced source has the actual A2 coordinates and inverse chord
maps R0E.60-\/-63; the point graph retains
\(\tau'/(\tau'-z^{2l+1}L)\), with \(L(0)=-1/4\). The good root
has the two full corrections
\(E_{2G'}-E_O\), with \(E_P(v_w,v_z)=\min(v_w,2lv_z)\), and the
actual ratio of their monomial generators. These are the complete
maps R0E.64-\/-70, including all strict-section units and their inverses.
Their full coefficient orders give \(2m/b\) for \(U_G\) at the root.
The normalized chord lift and its central-one factor yield the full
norm \(B'_{\rm sel}\), central value four; taking its prescribed root
retains the full finite character and gives
\(A_R=R^n\zeta_m^d\widetilde F^{\,n}D\),
\(A_R^m=B_R^n\), with \(R^m=\lambda\). The original selected
height factor is \(h=1+4z^m/27=z^mq\).

The full finite and root divisor equations for an automorphism are
\(\operatorname{ord}_0H=3j/b\),
\(\operatorname{ord}_\infty H=-2j/b\), and
\(\operatorname{ord}_{t_c}H=-k\). The last integer is its literal
selected \(B^k\) reindex. There are no other divisors. Their degree
gives \(j=bk,H=c(t^3/q)^k\). Consequently the full quotient map is
the physical scalar \(c\lambda^k\). These scalars extend through
every root and fan by the stated unit maps. Their kernel is trivial,
since an identity scalar deck would translate the source by
\(-kbG\) and \(G\) is not torsion. Thus
\(\operatorname{Aut}_{\rm hol}(X)=\mathbb C^*_{\rm physical}\).

For an endpoint the two independent period vectors \(e_1,\delta\)
force its complex-linear map to be \(eI\). Its suspension equation
retains an integral central difference and forces equality of the
two original scalar values. Write the resulting path winding as
\(J\in\mathbb Z\). The full annular reference R0E.72-\/-77 gives
the extra root factor \(z^{nJ}\); its integral suspension labels
have ratio \(e^{2\pi iJ J_\alpha}=1\). The divisor degree is now
\[
 3j/b-k+(-2mj/b-nJ)/m=0,
 \qquad j=bk+bnJ/m.                              \tag{R0R.17}
\]
Thus \(m\mid bJ\). All constants and reference cochains are retained
in the same ordinary/root comparison.

In the good row the finite map is \(S=tU_G\). On its good root,
\(f_{2,\rm old}=z^{-l}\phi\), where
\(\phi=(Y-1-z^lX/2)/(X-z^{2l}/4)\), and
\(B'_{2,1}(p)=\phi(p)\phi(p-G')\). The correct inverse is
\[
 U'_G(p,v)=(p+G',B'_{2,1}(p+G')^{-1}v).
 \tag{R0R.18}
\]
On the central cubic \(Y^2=X^3+1\), the two source maps have
\(X_+=2(1-Y)/X^2\) and \(X_-=2(1+Y)/X^2\).
Their actual line coefficients are respectively
\((Y+1)/(2X)=-1/X_+\) and \(-X\). Both compositions have coefficient
one at the inverse source point. Their three chronological products
have coefficients \(1/4\) and 4. This checks the inverse typing
without replacing its complete meromorphic Hom maps by those values.

The full root square is \(U_{G,\rm old}=z^{2l}U'_G\), read with
both corrections in R0R.16. Its deck relation is
\(U'_GD=\zeta_m^{-2l}DU'_G\). For \(J=l\) the actual forward
maps are
\[
 (R^{-l}tU_G)^n\quad\hbox{and}\quad(R^{-l}U'_G)^n.
 \tag{R0R.19}
\]
They intertwine the full finite generator with that at
\(R\mapsto\zeta_3R\), including \(\zeta_m^d\), since
\(-2ln\equiv ln\pmod {3l}\). The full annular square has product
\(z^{nl}t^nz^{2ln}=1\), using the original \(t=z^{-3l}\).
For the negative norm every factor is evaluated at its actual inverse
source point. The three chronological root phases give
\(\lambda^{-n}\); using \(S^3=qB_{\rm sel}^{-1}\), the composite
is the full quotient pair \(q^n,h^n\). In the reduced row that pair
itself is primitive at \(J=m\), since \(z^{nm}q^n=h^n\). Its
selected zero is carried by the literal \(B^n\) reindex and the
nonzero torus-unit map, with inverse, on every fan. Powers and inverses
give every multiple allowed by R0R.17.

\par

\medskip\noindent\textbf{R0R.4. Single-valued negative lifts and the completed endpoint group}

\par\nobreak

The selected inversion \((Z,V,q)\mapsto(Z^{-1},V^{-1},q)\) reverses
both full periods. Negating a triangle in a lattice-basis periodic fan
gives its opposite triangle with translated vertex \(-p-u-v\).
The actual chart maps R0E.90 exchange the two vertices, preserve
\(q\) and retain \(\kappa g,\alpha\eta\) and their reciprocals.
Thus this is an extending selected map for every stated fan.

A negative endpoint has a single-valued ordinary nonzero-line lift:
the entire finite III cover is proved in R0R.1, its \(O\)-section
vector kills the base meridian, and the other vanishing cycles lie
in \(\ker\psi\). Composing with the selected inversion gives the
actual positive selected-unit map and its inverse, so its continuation
has no additional selected monodromy. Constant integral deck joins
give the ordinary lift globally. On the normalized infinity root
fibre product, the same \(O\)-section vector in the actual point/O
frame R0E.65 or 68 gives a single-valued root lift. Both lifts reverse
the full suspension: \(\widetilde F B=B^{-1}\widetilde F\).

The complete annular reference around an original infinity circle
is \(AB^\nu\), with its actual retained integer \(\nu\). On a
root circle it is \(B^N\), where \(N=n+m\nu\). At both scalar
endpoints these are the same operators by the full reference
restoration. Choose one constant deck to equate the two lifts on
a cut annulus. Continuing that equality around the root circle gives
\[
 B^N\widetilde F B^{-N}=B^{2N}\widetilde F.
 \tag{R0R.20}
\]
Single-valuedness forces \(B^{2N}=1\), and strict contraction gives
\(N=0\). But \(N\equiv n\equiv\pm1\pmod m\), with \(m\ge2\),
so this is impossible for every retained reference integer. This
also treats the quadratic reduced cap. All negative endpoints are
excluded by their actual ordinary/root comparison.

The full winding group is consequently \(m\mathbb Z\) in the
reduced M6 row and \(l\mathbb Z\) in the good M2 row, with primitive
maps R0R.19 and the \(q^n,h^n\) pair just constructed. At a fixed
winding, the ratio of any two endpoints is exactly a physical scalar
automorphism. The proof retains both norm signs, every finite
character, every prescribed complete reference and every selected
lattice-basis fan. The executable peer independently derives the
Miller factors and both inverse-source compositions; its full kernel
and central-cochain equations are symbolic identities for all their
integer entries. These exact checks support the full geometric proof
above and the complete R0E source; they do not replace its extension,
compactness or single-valuedness arguments.

% End expanded verified cp36 input

\par\bigskip\noindent\textbf{Complete independent source, point-line, root and endpoint proof}\par\nobreak

% Expanded verified cp36 input: notes/part3_e6_iii_endpoint_fourfold_peer.tex
The source and its height are fixed literally as
\[
 E_0:\ y^2=x^3+tx+t^2,\qquad G=(0,t),\qquad
 t_c=-4/27,\qquad q=t+4/27.
 \tag{R0P.1}
\]
The completed objects retain the ordinary order-one, zero-norm III
attachment of K96 and either infinity cap
\((a,b;m)=(6,1;3l+2)\), \(l\ge0\), or
\((2,3;3l)\), \(l\ge1\), with \(n=\pm1\).
The genuine point line is \(M_a=\mathcal O([aG]-[O])\), the
suspension translates by \(-bG\), and all finite characters,
annular reference data, prescribed lattice-basis fans and admitted
strictly contracting scalars are retained. This independent proof
checks the complete R0E.1-\/-11 and R0R.1-\/-4 maps on that scope.

\par

\medskip\noindent\textbf{R0P.1. The whole finite source and its inverse graph maps}

\par\nobreak

The change
\(U=y-t-x/2\), \(V=y+t+x/2\) gives
\(UV=x^2(x-1/4)\), with inverse
\(t=(V-U-x)/2\), \(y=(U+V)/2\). At its A1 point the actual
\(x\)-chart of the resolution is
\[
 x=Y_s^2-T_s-T_s^2,\qquad
 (x,y,t)=(x,xY_s,xT_s).
 \tag{R0P.2}
\]
The fibre is \(xT_s=0\). The intersection algebra of its two
components is \(\mathbb C\{Y_s\}/(Y_s^2)\), so their intersection
number is two. Their total fibre has zero intersection with each
component, giving the matrix
\(\left(\begin{smallmatrix}-2&2\\2&-2\end{smallmatrix}\right)\)
and primitive kernel \((1,1)\).

The complete finite quartic map is
\[
 t=z^4,\qquad x=z^2X,\qquad y=z^3Y,
 \qquad Y^2=X^3+X+z^2,
 \quad D(z,X,Y)=(iz,-X,iY).
 \tag{R0P.3}
\]
The original equation pulls back with factor \(z^6\), and all three
original coordinate functions are invariant under this deck. On
\(X\ne0\), its map to R0P.2 is
\(T_s=z^2/X\), \(Y_s=zY/X\), including the two root points
\((X,Y)=(i,0),(-i,0)\). Substitution gives exactly R0P.2.

Translation by \(-G'\), \(G'=(0,z)\), has the full coefficients
\[
 L=(Y+z)/X,\qquad X_T=L^2-X,
 \qquad Y_T=L(X-X_T)-Y=z-LX_T.
 \tag{R0P.4}
\]
They preserve the whole equation. At the target point the slope for
\(+G'\) is \((Y_T-z)/X_T=-L\); its two coordinates return \(X,Y\).
This gives both actual source inverses before taking a central limit.

At \(O\), write \(X=U/\eta^2\), \(Y=-U/\eta^3\), where
\(U^2-U^3-\eta^4U-z^2\eta^6=0\) and \(U(0,0)=1\).
The exact pullback of R0P.4 is
\[
 X_T=\eta(\eta-2z)/U+2z^2\eta^4/U^2,
 \qquad
 Y_T=z+(1/\eta-z\eta^2/U)X_T.
 \tag{R0P.5}
\]
The first factor in \(X_T\) cancels the displayed \(1/\eta\).
The first jet of \(Y_T\) is \(\eta-z\), and its coefficient of
\(\eta\) is a unit. Thus \((Y_T,z)=(\eta,z)\) as analytic
ideals. At the other root centre set \(W=Y-z\) and
\(A=1-XW^2/(1+X^2)^2\). Reducing R0P.4 by the actual root equation
gives
\[
 -X_T/Y_T=
 \frac{W A(1+X^2)^2}{A(1+X^2)^3-zW^3}.
 \tag{R0P.6}
\]
Its denominator is a unit and its coefficient of \(Y-z\) is a
unit. Hence \((\eta_T,z)=(Y,z)\). The inverse translation supplies
the inverse ideal equalities. The universal property of the blowup
therefore gives morphisms and inverses on both entire projective
graphs, including all their exceptional directions.

The original-coordinate charts of these graphs are precisely
\(\eta=zp\), \(z=\eta v\), and \(Y=z\beta\), \(z=Yv\)
in R0E.17,20,21. In the last pair the source equations are
\(X=z^2\alpha\), \(\beta^2=1+\alpha+z^4\alpha^3\), and
\(X=Y^2C\), \(C(1+Y^4C^2)=1-v^2\), respectively. Together
with the \(X\)-unit chart R0P.2 these maps cover the whole compact
central root and its tangential source intersection.

\par

\medskip\noindent\textbf{R0P.2. Full Miller functions and genuine point-module cancellation}

\par\nobreak

Addition on R0P.1 gives
\[
 2G=(1/4,-t-1/8),\qquad
 3G=(8t+64t^2,3t+96t^2+512t^3).
 \tag{R0P.7}
\]
Write the second point as \((x_3,y_3)\), put
\(\mu_3=(3x_3^2+t)/(2y_3)\), and retain
\(x_6=\mu_3^2-2x_3\),
\(y_6=\mu_3(x_3-x_6)-y_3\). These are the actual doubled
point, not only its divisor. The entire functions are
\[
 \begin{gathered}
 f_2=\frac{y-t-x/2}{x-1/4},\qquad
 f_3=\frac{x(y-t-x/2)/4}{(y+t+x/2)/4-(1/4+2t)x},\\
 f_{2,3G}=\frac{y-y_3-\mu_3(x-x_3)}{x-x_6},\qquad
 f_6=f_3^2f_{2,3G}.
 \end{gathered}
 \tag{R0P.8}
\]
The actual chord formula for \(f_2g_{2G,G}\) reduces to the
displayed \(f_3\) by the original cubic equation. Its divisor and
that of \(f_6\) are the point-to-tensor divisors used in R0E.23.

Under R0P.3 these satisfy exactly
\(f_2=z^2\phi_2\), \(f_3=z^2\phi_3\),
\(f_{2,3G}=z^2\phi_{23}\), \(f_6=z^6\phi_6\), with the
entire rational functions R0E.25. All four \(\phi\)\textquotesingle s change sign
under the stated quartic deck. At its central curve they are
\(2X,X,6X,6X^3\). The full translated coefficient is therefore
\[
 t^{-3}B_{6,1}(p)=
 \phi_2(p)^6\frac{\phi_6(Tp)}{\phi_6(p)}.
 \tag{R0P.9}
\]
Its central Hom coefficient is 64, with its horizontal point and
zero modules retained. This identity follows before any such
restriction and retains the source argument of every factor.

On the identity component use \(x=p^{-2},y=-p^{-3}\); on the
other component use \(x=t\alpha,y=t\beta\), with
\(\alpha+t\alpha^3=\beta^2-1\). The actual source maps have
\(\beta'=p-1\) and \(p'=\beta-1\). Directly in the full
functions R0P.8 their leading coefficients are
\[
 \begin{array}{c|cc}
 & C_0&C_1\\\hline
 f_2&2/[p(p-2)]&2t(\beta-1)^2+O(t^2)\\
 f_3&p^{-2}&t(\beta-1)^3/(\beta-3)+O(t^2)\\
 f_{2,3G}&6/[p(p-6)]&6t(\beta-3)^2+O(t^2)\\
 f_6&6/[p^5(p-6)]&6t^3(\beta-1)^6+O(t^4).
 \end{array}
 \tag{R0P.10}
\]
Multiplication with those directed source maps gives
\(t^{-a/2}B_{a,1}=2^a(p-a)/p\) on \(C_0\), and
\(2^a(\beta-1)/(\beta-a-1)\) on \(C_1\), for \(a=2,6\).
The genuine frame is \(e_{a,i}=\sigma_a/r_{a,i}\), where
\(r_{a,0}=(p-a)/p\) and \(r_{a,1}=1\). Thus the directed
coefficient is the raw coefficient times
\(r_a(Tp)/r_a(p)\), giving \(2^a\) on each component. Its
chronological \(b=6/a\) product gives 64.

The two point divisors are retained globally by
\[
 \frac{\eta-z\tau_{a,\mathrm{old}}(z^4)}\eta
 =\frac{p-\tau_{a,\mathrm{old}}(z^4)}p
 =1-v\tau_{a,\mathrm{old}}(z^4),
 \tag{R0P.11}
\]
on the respective graph charts. Here
\(\tau_{2,\mathrm{old}}=2/(1+8t)\) and
\(\tau_{6,\mathrm{old}}=-x_6/y_6\), with respective limits
2 and 6. The equal exceptional orders cancel in this explicit
ratio, and their distinct strict directions remain. The reciprocal
ratio supplies the inverse module morphism. The other graph centre
and the tangency chart carry no point divisor. The two raw Hom
coefficients tend to \(2^a\) at the tangency; their vertical
orders on both components are exactly \(a/2\). Subtracting the
whole-fibre divisor removes those orders. The remaining horizontal
Hom divisor is canceled by R0P.11. A rational section and its
reciprocal with zero divisor in these smooth local rings are units.
This proves the entire finite extension, including every point
where a displayed meromorphic coefficient alone has a zero or pole.

\par

\medskip\noindent\textbf{R0P.3. The full period module and the primitive elliptic quotient}

\par\nobreak

The completed sphere topology is carried by the actual K96 disc
family and the F35/F46 integral source, common blowup and ruling
maps. Their primitive normal circles and the full root/reference
maps are the specified ones in R0E.5,8. They give
\(H^2(X;\mathbb Z)=H^4(X;\mathbb Z)=0\) for each prescribed
fan. A horizontal polar divisor would have positive area against
a flat Kähler form on a regular compact two-torus, while its
restricted divisor class is zero. Thus every polar divisor is
vertical. Removal of codimension-two indeterminacy and compactness
make the meromorphic function constant on each regular fibre.
At a special component, invariance under the branches of
\(t-t_0=u^e\) forces its Laurent exponents to be divisible by
\(e\). This proves the precise equality
\(\pi^*\mathbb C(t)=\mathbb C(X)\). An endpoint and its inverse
therefore give an actual Möbius map on the original base.

Keep the genuine point-theta coefficient and all four periods
\[
 V_{\rm pt}=f_{P^a}(Z)v_M,\qquad
 \gamma=(-bg,\ell),\quad e_1=(1,0),\quad
 e_2=(\tau,ag),\quad\delta=(0,1).
 \tag{R0P.12}
\]
The full \(\ell\) is the logarithm of R0E.37, including
\(\lambda q/t^3\), every Miller factor and the translated theta
ratio. The tensor comparison is the complete constant R0E.38,
with both reciprocal Cartier module arrows.

The point-representative change \(g'=g+\mu+\nu\tau\) multiplies
the line coefficient by
\((-Z/P^a)^{a\nu}Q^{-a\nu(a\nu+1)/2}\). The identity
\(\Theta(Q^{-N}x)=(-x)^NQ^{-N(N+1)/2}\Theta(x)\), with
\(N=a\nu\), proves this exact formula. Its inverse linear
matrix is \(\left(\begin{smallmatrix}1&0\\-a\nu&1\end{smallmatrix}\right)\).
Applying it to all three target periods gives R0E.40, including
\(c=\ell'-\ell+ab\nu(2g+\mu+\nu\tau)\). The modular
change R0E.41 sends \(p e_1+r e_2\) and \(q_0e_1+s e_2\)
to exactly the new elliptic periods, because \(ps-q_0r=1\).
Its theta quotient fixes the entire additive constant. These
maps give the full monodromy
\[
 M_\alpha=
 \begin{pmatrix}1&0&0\\ bv_\alpha&A_\alpha&0\\
 c_\alpha&av_\alpha^tJA_\alpha&1\end{pmatrix},
 \qquad J=\begin{pmatrix}0&-1\\1&0\end{pmatrix},
 \tag{R0P.13}
\]
with its complete central product R0E.44.

The Abel class is nonzero on this source. Clearing a rational
coboundary and multiplying by six puts a multiple of \(G\) in
the identity components and gives a single-valued Lie logarithm.
At the finite III point, \(\tau(2G')=2z/(1+8z^4)\), and
\(\omega_t=z^{-1}\omega_{\rm root}\). Dividing the root
logarithm by \(z\) gives the extending original Lie coefficient.
The quartic elliptic action has no invariant period. At the
selected I1 the invariant annular period extends in \(xy=q\).
At infinity the good cubic logarithm of \(3G'\) has leading
term \(-z^2/8\), and division by \(z^2\) gives the original
Lie frame; its order-three action again has no invariant period.
The full transition is \(u_0(q_\infty)=q_\infty^{-1}u_t(q_\infty^{-1})\).
A global holomorphic coefficient with that transition is zero.
The original infinity value of \(3G\) is \((64,512)\), whose
additive parameter is \(-1/8\); all nonzero integer multiples
remain nonzero. This proves the required contradiction and
non-torsion through the actual finite and infinity Lie charts.

The vector equation \((A_\alpha-I)u+bzv_\alpha=0\) and its
dual, using
\(w^tJ(A-I)=(w-Aw)^tJA\), now give exactly
\(\mathbb Z\delta\) and \(\mathbb Z\psi\) as the invariant
integral subgroups, with \(\psi(\gamma)=1\) and zero on the
other three generators. These are primitive in the displayed
lattice. An endpoint carries them to the corresponding pulled-back
groups before any base rigidity is used.

Its periodic derivative on the universal cover is constant, and
complex linearity gives the exact affine map
\((z,w)\mapsto(Az+z_0,Cz+\epsilon w+w_0)\), with
\(A\ne0\), \(\epsilon=\pm1\). Its inverse is R0E.46.
The full lattice map on \(\ker\psi\) is
\(\left(\begin{smallmatrix}p&q_0&0\\r&s&0\\u&v&\epsilon\end{smallmatrix}\right)\).
The complex area identity makes its elliptic determinant one.
Consequently its first coordinate and inverse descend to actual
inverse elliptic source isomorphisms, with the line constants and
integers retained in the second coordinate. Directly from the
original cubic,
\[
 j_0(t)=6912/(27t+4).
 \tag{R0P.14}
\]
Equality of this invariant under the proved elliptic map forces
the Möbius map to be \(h(t)=t\).

\par

\medskip\noindent\textbf{R0P.4. Global source sections and the entire automorphism group}

\par\nobreak

The elliptic monodromy centralizer is \(\mathbb QI\): its
order-three infinity matrix generates \(\mathbb Q(\zeta_3)\),
and the selected unipotent has a nonzero nilpotent part. A
nonscalar element of that quadratic field generates it, so the
simultaneous commutant is scalar. Applying the two Abel
coboundary equations to R0P.13 then gives exactly
\(eI+c\delta\psi\), with all central cochain entries retained.
Integral unimodularity gives \(e=\pm1\); complex linearity
eliminates the real rank-one remainder. On a single completed
object the negative sign gives sixteen transverse fixed points
on every regular two-torus. Their closed analytic equalizer has
a horizontal compact curve component. Its positive degree over
the base contradicts its pairing with the zero class
\(\pi^*H^2(\mathbb P^1)\). The automorphism has positive
linear part.

The actual ordinary nonzero-line cover has a nonzero line vector
over the \(O\)-section of each whole ordinary source disc, using
R0P.11 at point coincidences. This kills the base meridian.
The remaining vanishing cycles lie in \(\ker\psi\), so the
endpoint lift and its inverse extend there. At the selected
source, the integral orders in a chart \(q=z_0z_1z_2\) sum to
one. A literal \(h^rB^k\) reindex removes them, retaining all
four full units of R0E.50. The remaining map on each character
is multiplication by \(K(q)^uC(q)^v\), with reciprocal inverse,
as in R0E.52.

The normalized infinity fibre product has its actual root lift
because \(A_\lambda^m=B_\lambda^n\), \(n=\pm1\). The
commutator with the cap generator is \(B_\lambda^k\).
Its \(m\)-fold iteration is \(B_\lambda^{mk}=1\), which
strict contraction makes \(k=0\). The integer joins on the
one-punctured base with its single order-\(m\) orbifold point
have zero homomorphism to \(\mathbb Z\). Their actual deck
powers therefore join these source projections to a global
source section.

The rational surface R0P.2 has fibres IV*, III and I1, Euler
number 12 and rank-ten unimodular divisor lattice. The actual
III direction of \(G\) is \([x:U:V]=[0:0:1]\), off the
strict fibre plane \(V-U-x=0\); it meets the nonidentity
component once. At infinity the two section arms are the
distinct root points \((0,1)\) and \(O\), with the two
quotient rays specified in R0E.53. Thus the E6 and III
corrections are \(4/3\) and \(1/2\), and \(G\cdot O=0\).
The full ten-class matrix R0E.53 has determinant \(-1\).
These actual classes span the whole lattice; quotienting by
the fibre, zero section and seven components gives the free
section group \(\mathbb ZG\). In particular the source
section is \(jG\), with the actual primitive generator.

The full endpoint lift is \(H(t)U_G^j\), where
\(U_G(p,v)=(p+G,B_{a,1}(p+G)^{-1}v)\). Its finite orders
are \(-3/b\) on both III components by R0P.9-\/-11. The
infinity order is \(2m/b\) in its own root coordinate, and
the selected integer is the retained \(B^k\) reindex. The
three divisor orders of \(H\) are therefore
\(3j/b,-2j/b,-k\). Their sum gives \(j=bk\), and the
literal map is
\[
 H=c(t^3/q)^k,\qquad H U_G^{bk}=c B_{\rm sel}^{-k}.
 \tag{R0P.15}
\]
On the complete quotient this equals the physical scalar
\(c\lambda^k\). The reciprocal maps extend through the same
source, root, fan and reference charts. An identity scalar
deck would translate the source by \(-kbG\); non-torsion
forces \(k=0\). This proves precisely
\(\operatorname{Aut}_{\rm hol}(X)=\mathbb C^*_{\rm physical}\),
including its faithful action and connected scalar paths.

\par

\medskip\noindent\textbf{R0P.5. Both complete infinity-root and line morphisms}

\par\nobreak

The infinity coordinate is \(q_\infty=z^m\), so its two
original-coordinate maps are
\[
 \begin{array}{c|c|c}
 m&(t,x,y)&\text{root equation}\\\hline
 3l+2&(z^{-m},z^{-2(l+1)}X,z^{-3(l+1)}Y)
   &Y^2=X^3+z^{l+2}X+z^2\\
 3l&(z^{-3l},z^{-2l}X,z^{-3l}Y)
   &Y^2=X^3+z^lX+1.
 \end{array}
 \tag{R0P.16}
\]
These share the stated original \(t\) on the ordinary overlap
with R0P.3. Each root is used on its own original infinity
disc. The selected factor remains
\(h=1+4z^m/27=z^mq\).

For the reduced root put \(k=2l+1\) and
\(c=(1+z^lX)^{1/2}\), with its given unit branch. The
coordinates \(U=Y-zc,V=Y+zc\) have \(UV=X^3\), and
\(2zc=V-U\) has its actual analytic inverse for \(z\).
The three charts R0E.60 are regular and retain this source
map. In the middle chart \(X=ab,U=a^2b,V=ab^2\), write
\(v=z^l\). The whole R0E.61 chord has slope \(aK\),
\(X'=a(a-b)F\), \(Y'=-z-aKX'\), and
\(V'=-aX'H\). Our independent calculation uses the
relations \(c^2=1+vab\), \(2zc=ab(b-a)\) and
\((c')^2=1+vX'\) to prove these identities before taking
any central value. The coordinate map
\(a'=(b-a)F/H\), \(b'=-aH\) and the inverse chord
are therefore the actual two source morphisms, with
\(c,c+1,c'+1,K,F,H\) unit on these charts.

The vertical field in R0E.62 has
\(W(X)=2Y\) and \(W(Y)=3X^2+z^{l+2}\) on the fixed
\(z\) slice; these identities also hold in the independent
universal calculation. It is the extending invariant group
field. Since \(2m-1=3(2l+1)\), the small formal inverse
in R0E.63 has derivative \(1/m\) and supplies the stated
extending division translation with \(mR_{\rm div}=G'\).
The complete deck weights make every displayed unit
invariant. Its line graph is \((\tau',z^k)\), with all
rays \((j,1)\), \(1\le j\le k\). Adjacent determinants
are one, and both section equations have order \(j\) on
the corresponding ray. The two actual module ratios are
\(u/(u-L)\) and \(1/(1-vL)\), with \(L(0)=-1/4\).
Their inverses give the point-line identification carried by
the K96 common graph and commuting source/root maps.
The exact multiplier is \(h z^{-2}B'_6\); the full chord
gives \(B'_6/z^2=4\) at the centre, with its complete
horizontal modules retained.

For the good root set \(u=z^l\). The two marked charts
have quotient weights \((2,1)\), invariant equation
\(AB=C^3\), and rays \((2/3,1/3),(1/3,2/3)\).
The literal generator ratio gives
\(\beta^*N_2=\pi^*M_2+E_{2G'}-E_O\), with
\(E_P(v_w,v_z)=\min(v_w,2lv_z)\). The small translation
has \(w'=w\varepsilon(w,u)+u^2a(w,u)\), preserving
\((w,u^j)\) and \((u,w^j)\), \(j=1,2\), with their
inverse arrows. Thus both complete corrections are carried.
The whole pullback is
\(f_{2,\mathrm{old}}=z^{-l}\phi\), with the actual
\(\phi\) of R0E.69. The three central translated factors
are \((Y-1)/X,-(Y+1)/X,2/X\); their product is \(-2\),
and their squared norm is four. The selected multiplier
is the full \(hB'_{2,3}\).

For both rows the chord coefficient
\(g_{P_0,R}=\ell_{P_0,R}/(x-x(P_0+R))\), with
\(P_0=aG'\), has exactly the required horizontal Hom
divisor. Multiplication by the stated \(\tau(R)\) or
\(\tau(P_0)\) gives central coefficient one in its
actual module frames. Its two deck weights cancel. The
ratio \(\beta=B'_{\rm sel}/F^m\) is an invariant
extending base unit with central value four; its prescribed
analytic root retains the reference. Hence
\[
 \widetilde F=\beta^{1/m}F,\quad
 A_R=R^n\zeta_m^d\widetilde F^{\,n}D,\quad R^m=\lambda,
 \qquad A_R^m=B_R^n,\quad [A_R,B_R]=1
 \tag{R0P.17}
\]
are full source and line identities for both signs and every
finite character.

\par

\medskip\noindent\textbf{R0P.6. Full annular reference and the correct primitive inverses}

\par\nobreak

Use the actual root map \(I_{z_\alpha}\), transported
periods and reference integers from R0E.72-\/-74. With
\(s_\beta=s_\alpha+k\) and
\(f_\beta=D_*^kf_\alpha+\Pi_{\alpha\beta}\), the
entire residual is
\[
 h_\beta-D_*^kh_\alpha
 =-n s_\alpha\Pi_{\alpha\beta}
   +j_{\alpha\beta}\operatorname{Log}_\beta B_R
   +\Lambda_{\alpha\beta}.
 \tag{R0P.18}
\]
The written partition and Cauchy solution retain their
prescribed reference values and all integral columns.
Subtracting the target-sheet exponent
\(knf_\beta+kc_R\delta\) from the full exponent gives
exactly \(j_{\alpha\beta}\operatorname{Log}_\beta B_R+
\Lambda_{\alpha\beta}\). Thus
\(\Phi_\beta=A_R^kB_R^{j_{\alpha\beta}}\Phi_\alpha\),
with the displayed reciprocal inverse at its actual source
point. This verifies the directed annular square, including
the source-dependent period corrections.

The complete scalar family is R0E.76. At integral winding
\(J\) it has \(B_{\theta+J}=B_\theta\) and
\(\Phi_{\theta+J}=z^{nJ}\Phi_\theta\). The additional
sector scalar is \(e^{2\pi iJ J_\alpha}=1\) for its
retained integer, so no sector column is erased. The full
ordinary intertwiner and the independent period vectors
\(e_1,\delta\) force equality of the original scalar
values. The winding is therefore integral.

The three actual divisor orders for a positive endpoint are
R0E.79. Their degree gives \(j=bk+bnJ/m\), and hence
\(m\mid bJ\). For the good root the actual central
coefficient of \(B'_{2,1}\) is \(-X(p)\). Its inverse
is evaluated at \(p+G'\), giving precisely
\[
 U'_G(p,v)=\bigl(p+G',B'_{2,1}(p+G')^{-1}v\bigr),
 \qquad B'_{2,1}(p+G')^{-1}\big|_0
 =-1/X(p+G')=(Y+1)/(2X).
 \tag{R0P.19}
\]
The equality follows from the actual chord
\(X(p+G')=2(1-Y)/X^2=-2X/(Y+1)\) on
\(Y^2=X^3+1\). Our independent execution checks both
source inverses, and the three chronological plus-source
composites give the identity source and coefficient \(1/4\).
The complete root functions and their Cartier modules,
rather than these central values alone, supply the extension.

The literal ordinary/root square is
\(U_{G,\mathrm{old}}=z^{2l}U'_G\), with both graph
corrections from R0P.16. The entire deck relation is
\(U'_G D=\zeta_m^{-2l}D U'_G\). The normalized norm
lift commutes with it in the full constant-factor point
frame, retaining the base constants fixed by their
specified functions. Thus the primitive good forward maps
are
\[
 (R^{-l}tU_G)^n,\qquad (R^{-l}U'_G)^n,
 \tag{R0P.20}
\]
with \(-2ln\equiv ln\pmod {3l}\) and the exact
annular factor \(z^{nl}t^nz^{2ln}=1\).
For \(n=-1\) every reciprocal is evaluated at the
actual inverse source point.

Three chronological roots \(R,\zeta_3R,\zeta_3^2R\)
give scalar product \(\lambda^{-n}\). With
\(S=tU_G\), \(S^3=qB_{\rm sel}^{-1}\), their
complete endpoint product is
\[
 q^nB_\lambda^{-n}\quad\hbox{and}\quad
 h^nB_\lambda^{-n}.
 \tag{R0P.21}
\]
For M6 the primitive winding is \(m\), with the
canonical \(q^n/h^n\) quotient pair. Its selected
zero is carried by the actual \(B^n\) reindex and
the unit character map \((P^{bn},\eta^{-n})\) in
R0E.50, with reciprocal inverse on every cone.
This constructs every positive-part winding allowed
by the divisor equation, including all finite characters.

\par

\medskip\noindent\textbf{R0P.7. Negative lifts and the complete winding groups}

\par\nobreak

The selected inversion \((Z,V,q)\mapsto(Z^{-1},V^{-1},q)\)
reverses the two full periods. Negation sends a periodic
basis triangle to its other triangle with translated
vertex \(-p-u-v\). The full representative \(hBI\)
has the two chart maps R0E.90, with
\(\mathcal A=\kappa g\), \(\mathcal C=\alpha\eta\).
Their products preserve \(q\), their dense maps are
\(Z'=\mathcal A q/Z\), \(V'=\mathcal C q/V\), and
the two coordinate substitutions compose to the literal
identity. All coefficient functions remain evaluated at
the unchanged source height. Integral character changes
carry these arrows to every prescribed basis fan.

Composing a negative endpoint with that extending
inversion gives the positive selected unit map and
its inverse. It therefore supplies a single-valued
negative lift on the selected punctured disc.
On every other ordinary disc the actual point/O
frame gives the nonzero \(O\)-section vector killing
the meridian, and the vanishing cycles are in
\(\ker\psi\). The constant suspension joins give
one global ordinary lift. The normalized infinity
fibre product and its point/O Cartier frame give
the corresponding single-valued root lift fixing
\(z\). Both retain
\(\widetilde F B=B^{-1}\widetilde F\).

One original infinity circle has the full reference
word \(AB^\nu\), with the actual integer \(\nu\).
A root circle has word \(B^N\),
\(N=n+m\nu\). Both endpoints retain the same
integer by R0P.18 and the scalar restoration; the
extra \(z^{nJ}\) is single-valued on that root
circle. Equating the two lifts on a cut annulus,
with its one constant deck, and continuing around
the root gives
\[
 B^N\widetilde F B^{-N}=B^{2N}\widetilde F.
 \tag{R0P.22}
\]
Single-valuedness and strict contraction give
\(N=0\), while \(N\equiv\pm1\pmod m\) and
\(m\ge2\). This excludes every negative endpoint
with its actual annular integer, including the
quadratic reduced cap.

The complete winding groups are therefore
\(m\mathbb Z\) for \((6,1;3l+2)\) and
\(l\mathbb Z\) for \((2,3;3l)\), with the
primitive maps R0P.20-\/-21 and their chronological
inverses. At a fixed winding the full endpoint
ambiguity is the physical scalar group proved in
R0P.15. The physical return uses the inverse
endpoint. This proof assigns no numerical return.

\par

\medskip\noindent\textbf{R0P.8. Independent execution and audit scope}

\par\nobreak

The independent local executable is
\path{code/part3_fourfold/e6_iii_local_point_peer.py}.
Its 67 checks derive the full quartic source and
inverse, all four Miller functions, their component
limits and genuine-frame ratios, the two graph
ideal equations, tangency map and whole-fibre
intersection kernel. The independent global
executable is
\path{code/part3_fourfold/e6_iii_global_endpoint_peer.py}.
Its 1663 further checks retain the complete good
root source inverses and their threefold product,
the reduced chord in its full \(c,c'\) equations,
the invariant vertical field, every root exponent,
point graph, signed character, period coordinate,
target-sheet reference equation and endpoint word.
No author or root calculator is imported.

The full geometric read of R0E.1-\/-11 and R0R.1-\/-4
passes. The analytic extension, compact topology,
primitive invariant quotient, actual source section
and single-valued negative lifts are established
above using the retained explicit maps. The finite
integer executions supplement the universal
identities and those geometric proofs. The final
record locks all source and certificate bytes;
the two initial comparison-code stages are
preserved with their exact defining-equation
corrections. They changed no mathematical source.

% End expanded verified cp36 input

\par\bigskip\noindent\textbf{Independent proof of the full root, signed endpoint and annular comparisons}\par\nobreak

% Expanded verified cp36 input: notes/part3_e6_iii_endpoint_checks_peer.tex
\par

\medskip\noindent\textbf{R0CP.1. Exact scope and source dependency}

\par\nobreak

This is the independent checks audit of R0E.8-\/-10 at source SHA256
\texttt{21a7e21c\allowbreak 39ceedef\allowbreak 3abf4c14\allowbreak 8362ecd7\allowbreak a480d063\allowbreak fcc0c1c4\allowbreak 25fe6c2a\allowbreak 95c0c9ce}.
The original source is
\[
 y^2=x^3+tx+t^2,\qquad G=(0,t),\qquad
 t_c=-4/27,\qquad q=t+4/27.
\]
The genuine point line is \(M_a=\mathcal O([aG]-[O])\), and the
suspension translates by \(-bG\). Both entire cap progressions
\((a,b;m)=(6,1;3l+2)\), \(l\geq0\), and
\((2,3;3l)\), \(l\geq1\), retain \(n=\pm1\), the prescribed finite
character, the original selected-height fan and the full K96 reference.

The finite R0E.1-\/-4 source and Cartier maps have their separate
49-check audit. For the present read I also read all R0E.5-\/-7:
the whole point-theta period module, its full suspension coefficient,
the nonzero Abel class in the explicitly identified Lie line, both
primitive invariant lattice groups, the affine source map and inverse,
the rational source invariant, the centralizer, the actual cover lifts,
and the integral source-section lattice. The present numerical checks
do not replace those geometric arguments.

The stated negative and positive endpoint conclusions use those actual
cover lifts. At an ordinary III disc, a nonzero vector over the full
\(O\)-section disc kills the base meridian; the remaining generic-fibre
loops are subject to the actual vanishing-cycle relations in
\(\ker\psi\). The same argument uses the original point/O Cartier
modules on the normalized root disc. Thus the later annular comparison
uses single-valued lifts on both actual covers.

\par

\medskip\noindent\textbf{R0CP.2. Both original-coordinate root and line maps}

\par\nobreak

In the reduced row, the literal substitutions are
\[
 t=z^{-(3l+2)},\quad x=z^{-2(l+1)}X,\quad y=z^{-3(l+1)}Y,
 \qquad Y^2=X^3+z^{l+2}X+z^2.
\]
With \(c=(1+z^lX)^{1/2}\), the coordinates
\(U=Y-zc,V=Y+zc\) satisfy \(UV=X^3\). The three displayed A2
charts, the middle-chart chord and its inverse preserve the same
\(z\). The identity formal inverse in
\(R_{\rm div}=2G'-[m]^{-1}((2m-1)G')\) has derivative \(1/m\),
and \(mR_{\rm div}=G'\). The deck weights retain the invariant
\(c\) and the complete chord units.

The point/O graph has \(\tau'=z^{2l+1}\tau_0\) and
\(\tau'(6G')=z^{2l+1}L\), \(L(0)=-1/4\). Its actual coefficient
ratios are \(u/(u-L)\) and \((1-vL)^{-1}\) on the two charts.
The two sections have equal exceptional orders on each preceding
ray and distinct strict directions at the equality ray. These
are the directed Cartier maps in
\(\beta^*N_6=\pi^*M_6\). The complete selected Hom coefficient is
\(h z^{-2}B'_6\), with \(h=1+4z^m/27\); its central coefficient
after that displayed factor is four.

In the good row the full map is
\[
 t=z^{-3l},\quad x=z^{-2l}X,\quad y=z^{-3l}Y,
 \qquad Y^2=X^3+z^lX+1.
\]
The quotient and both section directions retain the monomial
correction
\(E_P(v_w,v_z)=\min(v_w,2l v_z)\), so the comparison is
\(\beta^*N_2=\pi^*M_2+E_{2G'}-E_O\).
The small formal translation has the actual form
\(w\mapsto w\varepsilon+u^2a\), \(u=z^l\), with inverse;
it preserves all four ideals \((w,u^j),(u,w^j)\), \(j=1,2\).
The Miller factor remains
\[
 f_{2,\mathrm{old}}=z^{-l}\phi,\qquad
 \phi=\frac{Y-1-z^lX/2}{X-z^{2l}/4}.
\]
Its three central translates have product \(-2\); the full squared
norm is four. The displayed quotient monomials remain in the
ordinary/root comparison.

In each row the normalized chord-over-vertical map has exactly the
Hom divisor for translation by \(-R_{\rm div}\). Multiplication by
the stated formal \(\tau\)-factor cancels its pole and gives central
coefficient one. Its opposite deck weights cancel before taking the
reference-fixed root of \(B'_{\rm sel}/F^m\). Consequently the full
maps satisfy
\[
 \widetilde F^m=B'_{\rm sel},\qquad
 A_R=R^n\zeta_m^d\widetilde F^{\,n}D,\qquad
 A_R^m=B_\lambda^n,
\]
with every base unit and finite character present. The finite III
quartic charts and these infinity charts are compared only through
their displayed original \(t,x,y\) maps on the ordinary overlap.

\par

\medskip\noindent\textbf{R0CP.3. The actual inverse coefficient and its chronological norm}

\par\nobreak

On the central good cubic \(Y^2=X^3+1\), the chord for translation
by \(+G'=(0,1)\) gives
\[
 X_1=-\frac{2X}{Y+1},\qquad Y_1=\frac{Y-3}{Y+1}.
 \tag{R0CP.1}
\]
The full forward \(-G'\) Hom coefficient has central expression
\(B'_{2,1}(p)=-X(p)\) in the canonical meromorphic point frame.
Its inverse is therefore
\[
 U'_G(p,v)=
 \left(p+G',\frac{Y+1}{2X}v\right),
 \qquad
 \frac{Y+1}{2X}=-\frac1{X(p+G')}.
 \tag{R0CP.2}
\]
The original point/O modules carry the zeros and poles of this
meromorphic expression and its reciprocal. Their Hom maps have
zero vertical order and extend on the actual root line.

The next source point is
\(X_2=2X/(Y-1),Y_2=-(Y+3)/(Y-1)\), and the third is the original
point. The three coefficients of the inverse arrow, evaluated at
their chronological source points, are
\[
 \frac{Y+1}{2X},\qquad -\frac{Y-1}{2X},\qquad-\frac1X.
\]
Their product is \((Y^2-1)/(4X^3)=1/4\). This gives a pointwise
check of the inverse direction and its central norm. The earlier
wording that identified this inverse coefficient with \(-X(p)\)
was located in the audit and preserved by the author before repair.
The complete ordinary/root square is still
\(U_{G,\mathrm{old}}=z^{2l}U'_G\), with its actual correction
monomials; the corrected source-point evaluation is now literal.

\par

\medskip\noindent\textbf{R0CP.4. Full annular reference, scalar family and signed endpoints}

\par\nobreak

The sector construction retains the complete source displacement,
physical scalar and root change in
\[
 \Phi_\alpha=\operatorname{Exp}
 (ns_\alpha f_\alpha+s_\alpha c_\lambda\delta+h_\alpha)I_{z_\alpha}.
\]
Its cochain equation carries both the full period \(\Pi\) and every
integer \(j,\Lambda\). The reference-zero Cauchy solve is performed
on that residual. Subtraction of the actual target-sheet logarithm
then gives \(\Phi_\beta=A_\lambda^kB_\lambda^j\Phi_\alpha\).
The inverse uses the actual inverse source point and reciprocal
Cartier maps.

Keeping the integer representatives \(J_\alpha\), the scalar family
is exactly
\[
 \Phi_{\theta,\alpha}=e^{2\pi in\theta s_\alpha/m}
 B_\theta^{J_\alpha}B_*^{-J_\alpha}\Phi_{*,\alpha}.
\]
The overlap exponent is \(n\theta k/m+\theta j\). At integer
winding \(J\), the extra exponent is
\(nJs_\alpha/m+JJ_\alpha\); the second term is an integer after
it has been carried through the square. Thus the full residual is
\(z^{nJ}\).

For a positive endpoint \(H(t)U_G^j\), the actual finite, selected
and root divisors give
\[
 \operatorname{ord}_0H=3j/b,\quad
 \operatorname{ord}_{t_c}H=-k,\quad
 m\operatorname{ord}_\infty H+2mj/b+nJ=0.
\]
Their degree gives \(j=bk+bnJ/m\) and \(m\mid bJ\).
This is the integral restriction on the full rational base
coefficient, with the selected reindex \(B^k\) retained.

For M2, \(S=tU_G\) has the entire selected coefficient of R0E.81,
an extending unit and reciprocal at unchanged \(q\), on every
prescribed cone. On the good root
\(U'_GD=\zeta_m^{-2l}DU'_G\). The genuine point-extension
translations commute with the full norm lift; their base constants
remain in the constant-factor point frame. Hence, for either sign,
\[
 (U'_G)^nA_R=A_{\zeta_3R}(U'_G)^n.
\]
The primitive forward maps are
\((R^{-l}S)^n\) and \((R^{-l}U'_G)^n\). Their annular square
cancels precisely
\(z^{nl}t^nz^{2ln}=1\), using the original \(t=z^{-3l}\).
The three successive root basepoints contribute the complete scalar
\(\lambda^{-n}\), so the chronological products are
\(q^nB_\lambda^{-n}\) and \(h^nB_\lambda^{-n}\).

For M6 the primitive forward \(q^n/h^n\) closing uses
\(z^{nm}q^n=h^n\). At the selected zero its literal
\(B_\lambda^n\) reindex leaves the source/line unit pair
\((P^{bn},\eta^{-n})\); each character and reciprocal extends
at unchanged height. Negative winding and negative primitive norm
use the corresponding chronological inverse maps.

\par

\medskip\noindent\textbf{R0CP.5. Complete negative continuation and audit conclusion}

\par\nobreak

The selected inversion \((Z,V,q)\mapsto(Z^{-1},V^{-1},q)\)
conjugates both full periods to their inverses. The two explicit
\(hBI\) vertex maps preserve \(q\), interchange the charts and
are literal inverses, with \(\mathcal A=\kappa g\) and
\(\mathcal C=\alpha\eta\) retained. Composing a negative lift
with this map gives the positive integral-height unit construction;
composing back supplies its single-valued selected lift.

On the normalized root the preserved kernel is exactly the
\(B_\lambda\) cyclic cover, since \(n=\pm1\). The actual Cartier
frames over the \(O\)-section give its single-valued root lift.
The entire original-base turn is \(A_\lambda B_\lambda^\nu\),
and its full root turn is \(B_\lambda^N\),
\(N=n+m\nu\). The integer \(\nu\) remains unrestricted.

Continuing the directed comparison
\(\widetilde F_{\rm cap}=\Phi_t\widetilde F_{\rm ord}\Phi_s^{-1}\)
around that root circle gives
\(B_\lambda^N\widetilde F_{\rm cap}B_\lambda^{-N}
=B_\lambda^{2N}\widetilde F_{\rm cap}\).
Both lifts are single valued and the contraction has infinite
order, so this requires \(N=0\), whereas
\(n+m\nu\) cannot vanish for \(n=\pm1,m\geq2\).
Every analytic coefficient is already in the displayed operators;
this contradiction does not remove a reference integer.

The complete R0E.8-\/-10 geometric read passes after the preserved
inverse-coefficient repair. The independent certificate checks the
actual central source/inverse, both vertex inverses, full root and
finite-character exponents, arbitrary sector representatives,
the three-divisor equation and the entire negative annular word.
Its finite checks accompany the explicit extension, covering and
cochain arguments above. This audit assigns no numerical return.

% End expanded verified cp36 input

\end{proof}


\subsection{Complete scalar returns on the fixed E6 source locus}
